% !TeX program = pdflatex
% !TeX encoding = utf-8
% ------------------------------------------------------------------
%  2026 ShuWei Cup, Problem B
%  Fast UAV Localization and Scheduling with Multiple Cooperative
%  Mobile Observation Vehicles
%
%  English LaTeX version built on the official swmcmthesis template.
%  Compile with pdfLaTeX (two passes, so the TOC and page refs settle):
%      pdflatex main
%      pdflatex main
% ------------------------------------------------------------------
\documentclass{swmcmthesis}

% ---- page geometry: A4, compact margins to respect the 25-page limit
\geometry{top=20mm,bottom=20mm,left=20mm,right=20mm,headsep=6mm}
\setlength{\headheight}{18pt}

% ---- body typography
\linespread{1.3}
\setlength{\parindent}{2em}
\setlength{\parskip}{0pt}
\setcounter{tocdepth}{2}

% ---- chapter-based figure/table numbering (Figure 2-1, Table 4-3, ...)
\makeatletter
\@addtoreset{figure}{section}
\@addtoreset{table}{section}
\makeatother
\renewcommand{\thefigure}{\arabic{section}-\arabic{figure}}
\renewcommand{\thetable}{\arabic{section}-\arabic{table}}

% ---- notation shortcuts
\newcommand{\rmin}{r_{\min}}
\newcommand{\wrap}{\operatorname{wrap180}}
\newcommand{\distpath}{\operatorname{dist}_{\mathrm{path}}}
\newcommand{\diag}{\operatorname{diag}}
\newcommand{\trace}{\operatorname{trace}}
\newcommand{\argmax}{\operatorname*{arg\,max}}
\newcommand{\argmin}{\operatorname*{arg\,min}}
\newcommand{\figwidth}{10.4cm}
\newcommand{\bsl}{/\allowbreak}

% ---- unbreakable table body --------------------------------------
% The template's \tabular forces \arraystretch=1.3, which makes the
% taller tables overflow a single page.  tabularnb restores the plain
% tabular so the per-table \arraystretch is honoured, and, being an
% ordinary tabular, it never splits across pages.
\makeatletter
\newenvironment{tabularnb}{\swmcm@oldtabular}{\swmcm@endoldtabular}
\makeatother

% ---- summary sheet -------------------------------------------------
\title{Fast UAV Localization and Scheduling with Multiple Cooperative Mobile Observation Vehicles}

% Team control number (ASCII; the template is compiled with pdfLaTeX).
\baominghao{SWBMCM2611718}

\zubie{B}
\date{2026}

\abstract{%
{\linespread{1.15}\selectfont\footnotesize
Low-altitude small UAVs are reported only by coarse alerts---approximate azimuth, horizontal range and altitude---whose nominal values do not amount to a precise localization. Following the chain ``geometric feasibility---intersection accuracy---error propagation---disturbance response'', this paper builds four models---observable-region determination, two-station intersection accuracy, pointing-error propagation and event-triggered rescheduling---and follows the route of \textbf{certifying the optimal value in closed form, enumerating the optimal solution deterministically, and using numerical experiments only for verification}.

\textbf{For Problem~1} (observable region and optimal station for one target), the three constraints---vehicle geometry, gimbal pitch and line-of-sight (LOS) occlusion---are intersected, and the reciprocal of the horizontal distance is used as the quality measure, \textbf{giving an observable-region model and a quality-optimization model}. \textbf{The minimum horizontal distance is about 7.58~m} with the elevation at \textbf{the $80^\circ$ pitch limit}, so the optimum is attained with the pitch constraint active and is independent of the building layout, and the optimal set is \textbf{a continuous arc of 283.85$^\circ$}. A representative station $(88.88,76.29)$~m is then selected by the three-level rule (maximum safety-margin band, shortest reachable point within the band, coordinate lexicographic order). The observable area is 5610~m$^2$ (\textbf{56.1\% of the site}), and the grid recomputation differs from the closed-form bound by only $0.0007$~m.

\textbf{For Problem~2} (two-station intersection accuracy and multi-alert scheduling), the error propagation relation is derived from the LOS equations, the geometric dilution of precision $G$ characterizes the accuracy, and the mean inequality proves the lower bound, \textbf{giving a two-station accuracy model and a six-slot five-vehicle scheduling model}. Under the problem-specified $1^\circ$ azimuth error, \textbf{the accuracy bound is 5.8188~m}, attained when the two stations are equidistant and the lines of sight are orthogonal, with DRMS $0.1016$~m and a $95\%$ radius of $0.2486$~m. \textbf{The time-optimal schedule starts the latest alert at 10.355~s}, \textbf{50.6\% shorter than the nearest-assignment baseline of 20.968~s}, \textbf{but it is not deployable as it stands}: its first alert has an intersection angle of only $19.34^\circ$ and a DRMS of $3.084$~m, about $30$ times the accuracy bound. \textbf{Imposing the paper's $[30^\circ,150^\circ]$ operating window} gives an \textbf{accuracy-qualified plan whose latest start is 12.983~s} ($+25.4\%$), with a first-alert DRMS of $2.036$~m; \textbf{this plan is recommended for deployment}. For the scheduling optimum, a reuse-aware relaxation certifies a lower bound of $7.391$~s, so the $10.355$~s plan is within $28.6\%$ of optimal. Timeliness and accuracy are thus not interchangeable.

\textbf{For Problem~3} (localization and pointing errors versus horizontal azimuth accuracy), the localization error is projected onto the cross-range direction and linearized for small angles, and the three error sources are combined into a normal distribution, \textbf{giving a pointing-error propagation model}. \textbf{The maximum reliable observation distance is 66.92~m} (a conservative 66.84~m from high-accuracy quadrature), and the dominant error source switches at a vehicle--target distance of about $19.9$~m. At $80$~m the probability of meeting the criterion falls to $0.915$, so the distance-error coefficient must be reduced by about $16.3\%$; localization then contributes 88.1\% of the variance.

\textbf{For Problem~4} (plan failure under vehicle failure, communication outage and new alerts), event-triggered snapshots are combined with rescheduling, \textbf{giving a rescheduling model under disturbance}. Among 74 scenarios \textbf{a keep-original no-reaction baseline loses completion in 67}; a \textbf{simple repair} that only inserts idle vehicles and replaces the failed one by the nearest spare \textbf{already restores the completion ratio in all 67} (7 ties, no degradation), so the completion gain comes from reacting at all. \textbf{Re-optimization leaves the completion ratio unchanged but lowers the latest start in 61 of 74 scenarios} (mean $4.35$~s, maximum $13.10$~s; 13 ties, no degradation), e.g., a single vehicle failure goes from $60$~s (keep) to $21.01$~s (simple repair) and $11.17$~s (re-optimization). The mechanism is that \textbf{structural loss is converted into temporal loss}, and re-optimization compresses the temporal part. This $\eta$ accounting follows the problem's started-means-served rule; under a $\ge4$~s completion criterion the $\eta=1$ count falls to $66/74$, but the repair/re-optimization attribution is unchanged.

\textbf{Sensitivity analysis} shows that the pitch limit, the angular rate, the vehicle--target distance and the communication recovery time are the core parameters of the four problems: a 1~m rise in target height increases the minimum horizontal distance by 0.176~m, a $1^\circ$ relaxation of the pitch limit reduces it by 0.774~m, and the elasticity of the angular rate is about four times that of the travel speed; no qualitative conclusion changes under reasonable perturbations. The paper thus gives provable results for station selection and accuracy bounds, and reliable plans for scheduling and disturbances in the sense of structural exactness and station-wise local optimality, providing a quantitative basis for station placement and cooperative scheduling in low-altitude surveillance.%
\par}
}

\keywords{UAV localization; two-station intersection; geometric dilution of precision; observable region; vehicle scheduling; event-triggered rescheduling}

\begin{document}

\maketitle
\begingroup
\linespread{1.0}\selectfont
\tableofcontents
\endgroup
\newpage

% ---------------- body: pages 1..25 -------------------------------
\setcounter{page}{1}
\rhead{\small Page \thepage\ of 25}

% =====================================================================
% 1. Introduction
% =====================================================================
\section{Introduction}
\label{sec:introduction}

\subsection{Background}

When a low-altitude small UAV enters the surveillance airspace, the reconnaissance system can only provide a coarse alert---an approximate azimuth, horizontal distance, and altitude---which does not give a precise position. The nominal alert values are therefore used only as planning data (Assumption~A5): they fix the target position for feasibility checks, station selection, and scheduling, but are neither a measurement nor a localization result. The precise horizontal position is obtained, as Problem~2 prescribes, by two-station synchronous direction finding: a single vehicle yields only one bearing line, while two vehicles observing synchronously from different directions intersect two bearing lines. The target altitude does not enter the intersection and enters only through elevation constraints, building occlusion, and pointing errors. The accuracy analysis of Problem~2 is thus a conditional analysis of this bearing-intersection process, whose only stochastic input is the specified azimuth error. A baseline that is too short, or an intersection angle that is too small or too large, amplifies the error; station geometry is therefore the key to accuracy.

\subsection{Problems to Be Solved}

\begin{enumerate}[label=(\arabic*)]
\item \textbf{Single-target observable region and optimal station.} For T1, determine and plot the observable region under the vehicle-geometry, LOS-occlusion and gimbal-pitch constraints. Using $Q=1/r$, find the optimal station, distance and elevation, and state the optimality criterion and tie-breaking rules.

\item \textbf{Two-station intersection accuracy and multi-target assignment.} Under azimuth errors $\mathcal{N}(0,(1^\circ)^2)$, relate accuracy to baseline, intersection angle and range, and define an accuracy index. Under a 30~s travel limit, select two of five vehicles for T2 and analyze degeneracy; then minimize the latest observation start for three alerts on six slots shared by five vehicles, output the timeline, and define a nearest-assignment baseline.

\item \textbf{Localization and pointing error propagation.} With $\sigma(s)=0.3+0.02s$~m, gimbal bias $0.8^\circ$, control error $0.6^\circ$ and a 40~m range, build the $\Delta\psi$ propagation model and find the maximum travel distance satisfying $P(|\Delta\psi|\le5^\circ)\ge0.95$. Derive the $(a,b)$ constraints at 80~m, justify the approximation, and compare the three error sources.

\item \textbf{Rescheduling under disturbances.} Taking the 2(2) plan as the baseline, at $t=5$~s consider a new alert ($\theta\in[0^\circ,360^\circ)$, $\rho\in[30,50]$~m, $H=50$~m), one vehicle failure, an outage $[5,t_{\mathrm{rec}}]$ ($t_{\mathrm{rec}}\in(5,30]$) and their combinations. Build a rescheduling model, lock started observations and allow redispatch, design representative scenarios, and report $(\eta,t_{\mathrm{last}})$ within $[0,60]$~s.
\end{enumerate}

% =====================================================================
% 2. Problem Analysis
% =====================================================================
\section{Problem Analysis}
\label{sec:analysis}

The four problems progress along the chain ``geometric feasibility---intersection accuracy---error propagation---disturbance response'': Q1 answers where observation is possible and where it is best; Q2 extends single-station geometry to two-station intersection and multi-alert scheduling; Q3 advances deterministic geometry to error probability; Q4 advances static scheduling to disturbance response. The feasible region, accuracy conclusions, and timeline of earlier problems form the input baselines of later ones. Figure~\ref{fig:framework} shows the input--output relations and the modeling chain.

\begin{figure}[htbp]
  \centering
  \includegraphics[width=\figwidth]{fig0_framework.png}
  \caption{Framework of the four problems and the overall modeling chain}
  \label{fig:framework}
\end{figure}

\subsection{Geometry and Parameters}

The site is a $100~\mathrm{m}\times100~\mathrm{m}$ plane with the southwest corner as the origin, the $x$-axis east and the $y$-axis north; the dispatch center is $O=(50,50)$. The site contains five axis-aligned cuboid buildings and five standby observation vehicles; their geometry and initial positions are listed in Table~\ref{tab:2-1}.

\begin{table}[H]
\centering
\caption{Building geometry and initial positions of the observation vehicles}
\label{tab:2-1}
{\small
\renewcommand{\arraystretch}{1.08}
\begin{tabularnb}{@{}p{20mm}p{12mm}p{40mm}p{42mm}p{18mm}@{}}
  \toprule
  Category & ID & Planar range (m) & Initial position $(x,y)$ (m) & Height (m) \\
  \midrule
  Building & B1 & $[22,34]\times[16,26]$ & --- & 20 \\
  Building & B2 & $[63,79]\times[59,71]$ & --- & 28 \\
  Building & B3 & $[10,24]\times[65,76]$ & --- & 12 \\
  Building & B4 & $[72,86]\times[22,34]$ & --- & 24 \\
  Building & B5 & $[41,51]\times[38,46]$ & --- & 32 \\
  Vehicle & V1 & --- & $(15,14)$ & --- \\
  Vehicle & V2 & --- & $(88,10)$ & --- \\
  Vehicle & V3 & --- & $(14,86)$ & --- \\
  Vehicle & V4 & --- & $(88,68)$ & --- \\
  Vehicle & V5 & --- & $(50,93)$ & --- \\
  \bottomrule
\end{tabularnb}
}
\par\smallskip
\begin{minipage}{\linewidth}
\noindent\footnotesize Note: All buildings are axis-aligned cuboids standing on $z=0$; a vehicle is an equivalent disk of radius $1.5$~m, and the tabulated position is the initial coordinate of its center (the ground projection of the camera). The three baseline alerts are given in Section~2.1; the complete parameter table is provided in \texttt{params\_table.md} in the supporting materials.
\end{minipage}
\end{table}

A vehicle is modeled as a disk of radius $1.5$~m, travels at a constant $3$~m/s, carries a camera $2$~m above the ground, and has a gimbal pitch range of $[-10^\circ,80^\circ]$ with horizontal and pitch angular rates of $30^\circ$/s; the slew time is the longer axis. During travel the gimbal remains level facing east, and it may turn toward the target only after arrival. A task consists of travel, slew, synchronized $5$~s observation by two vehicles, and reset; all phases are counted on the timeline. A line of sight tangent to a building is judged occluded. An alert is given as $(\theta,R,H)$, and the target's horizontal position is $(50+R\cos\theta,50+R\sin\theta)$; the baseline alerts are T1$(30^\circ,40~\mathrm{m},45~\mathrm{m})$, T2$(150^\circ,45~\mathrm{m},35~\mathrm{m})$, and T3$(285^\circ,38~\mathrm{m},55~\mathrm{m})$ (see~\cite{swmcm}); Problem~4 adds T4$(\theta\in[0^\circ,360^\circ),R\in[30,50]~\mathrm{m},H=50~\mathrm{m})$. These nominal alert values are planning data for all geometry, feasibility, and scheduling computations (Assumption~A5); every accuracy conclusion is conditional on the bearing and pointing error models of Problems~2 and 3.

\subsection{Analysis of Problem 1}

Problem~1 asks where observation is possible, where it is best, how ``optimal'' is defined, and how ties are broken; the observable region is the feasible set, and the optimal station is its optimizer. The three constraints are heterogeneous and the optimizer may be a one-dimensional continuum, so we use an analytics-first route with numerics only as evidence.

\subsection{Analysis of Problem 2}

Part~2(1) requires a two-station accuracy model, an accuracy index, station selection for T2 under the 30~s limit, and a degeneracy discussion; Part~2(2) minimizes the latest start with three alerts, six slots and five vehicles and outputs the timeline against a nearest-assignment baseline. The two parts share geometry but not objectives, and 2(1) is a conditional accuracy analysis of two-station bearing intersection.

\subsection{Analysis of Problem 3}

Problem~3 requires a propagation model for $\Delta\psi$, $s_{\max}$, the $(a,b)$ constraints at 80~m, justification of the approximation, comparison of the three error sources, and deployment implications. The difficulties are the different dimensions of the sources, the nonlinear $\operatorname{atan2}$, the two-sided criterion, and the monotonicity needed for the $(a,b)$ constraints.

\subsection{Analysis of Problem 4}

Problem~4 takes the time-optimal 10.355~s plan as the baseline and considers, at $t=5$~s, new alerts, one vehicle failure, communication outages and their combinations, evaluated by $(\eta,t_{\mathrm{last}})$. Disturbances occur mid-execution (locking, outage-triggered rescheduling, continuous new-alert parameters), so the comparison includes a simple-repair baseline to separate reaction from re-optimization.

\subsection{Algorithm Routes and Selection}

Table~\ref{tab:2-2} gives a compact route summary; full details are in Appendix~B.1--B.2.

\begin{table}[htbp]
\centering
\caption{Compact algorithm routes for the four problems}
\label{tab:2-2}
{\small
\renewcommand{\arraystretch}{1.08}
\begin{tabularnb}{@{}p{12mm}p{64mm}p{74mm}@{}}
  \toprule
  Problem & Method & Why \\
  \midrule
  Q1 & Closed-form $r\ge\rmin$ + $\rmin$-circle sampling + 3-D slab/visibility graph + lexicographic tie-break &
  $Q=1/r$ is strictly monotone, so the optimum is certified in closed form; the optimal set is a one-dimensional arc \\
  \midrule
  Q2(1) & Covariance propagation $\Sigma=A^{-1}DA^{-T}$ + QM--AM bound + $C(5,2)$ pair enumeration + polar candidates/Nelder--Mead &
  The equality condition reduces optimality to a $90^\circ$ pairing test on the $\rmin$ circle; only 10 pairs exist \\
  \midrule
  Q2(2) & 1,000 structures / 5,880 reuse orders + fixed-point timeline + coordinate descent/Nelder--Mead; reuse-aware LB &
  5 vehicles and 6 slots permit exhaustive deterministic enumeration; stations are locally refined; certified gap $28.6\%$ \\
  \midrule
  Q3 & Closed-form first-order propagation + two-sided quantile + $256\times256$ Gauss--Hermite quadrature + MC verification &
  $\sigma_\psi(s)$ is monotone, so the $P$ criterion is invertible; quadrature has no sampling error \\
  \midrule
  Q4 & Event-triggered snapshot + reuse of the Q2(2) solver + 74-scenario enumeration with keep/simple-repair/re-optimization comparison &
  The post-disturbance scale matches Q2(2), so the same solver is reused and all scenarios are comparable \\
  \bottomrule
\end{tabularnb}
}
\end{table}

% =====================================================================
% 3. Symbols and Assumptions
% =====================================================================
\section{Symbols and Assumptions}
\label{sec:symbols}

\subsection{Symbols}

\begin{table}[htbp]
\centering
\caption{Main symbols}
\label{tab:3-1}
{\small
\renewcommand{\arraystretch}{1.08}
\begin{tabularnb}{@{}p{38mm}p{92mm}p{30mm}@{}}
  \toprule
  Symbol & Meaning & Unit / value \\
  \midrule
  $O$ & Dispatch center & $(50,50)$~m \\
  $B_i$ & Building $i$ ($i=1,\dots,5$), planar range $[x_i^0,x_i^1]\times[y_i^0,y_i^1]$, height $h_i$ & m \\
  $T_a$; $\theta_a,\rho_a,H_a$; $T'_a$ & Alert $a$ and its target; azimuth (counterclockwise from east), horizontal distance, height; horizontal projection & $^\circ$, m, m; m \\
  $h_{\mathrm{cam}}$; $\Delta z_a$ & Camera height above ground; camera--target height difference $\Delta z_a=H_a-h_{\mathrm{cam}}$ & 2~m; m \\
  $p=(x,y)$; $r(p)$; $\varepsilon(p)$ & Center of the observation station; horizontal distance from camera to target projection; LOS elevation $\arctan(\Delta z_a/r)$ & m; m; $^\circ$ \\
  $[\varepsilon_{\min},\varepsilon_{\max}]$ & Feasible gimbal pitch range & $[-10^\circ,80^\circ]$ \\
  C1, C2, C3; $\Omega_a$ & Vehicle geometry, gimbal pitch, and LOS occlusion constraints; observable region $\Omega_a=\mathrm{C1}\cap\mathrm{C2}\cap\mathrm{C3}$ & ---; m$^2$ \\
  $Q(p)$; $p^*,r^*,\varepsilon^*,Q^*$ & Observation quality $1/r$; optimal station, horizontal distance, elevation, and quality in Problem~1 & m$^{-1}$; m, m, $^\circ$, m$^{-1}$ \\
  $m(p)$; $\varphi$ & Geometric safety margin; polar angle parameter of the optimal arc & m; $^\circ$ \\
  $V_j$; $V_j^0$ & Observation vehicle $j$ and its initial position ($j=1,\dots,5$) & m \\
  $R_{\mathrm{veh}}$; $v$; $\omega$ & Equivalent vehicle disk radius; travel speed; maximum gimbal single-axis angular rate & 1.5~m; 3~m/s; $30^\circ$/s \\
  $\alpha$ & Azimuth from station to target (pan command angle) & $^\circ$ \\
  $t_{\mathrm{tr}},t_{\mathrm{sl}},t_{\mathrm{obs}},t_{\mathrm{rs}}$ & Travel, slew, synchronized observation, and reset times & s ($t_{\mathrm{obs}}=5$) \\
  $d_i$; $d$ & Horizontal distance from station $i$ to target; single-station shorthand in Problem~3 (nominal 40~m) & m \\
  $\sigma_\alpha$; $\gamma$; $\beta$ & Single-station azimuth standard deviation; intersection angle of the two LOS; two-station baseline length & $1^\circ$ (rad in computation); $^\circ$; m \\
  $\Sigma$; $G$ & Covariance of the intersection-point position error; geometric dilution of precision $\sqrt{(d_1^2+d_2^2)/2}/|\sin\gamma|$ & m$^2$; m \\
  DRMS; $R_{95}$ & Root-mean-square position error $\sigma_\alpha G$; 95\% error-circle radius $2.4477\,\mathrm{DRMS}$ & m; m \\
  $\pi_j$; $p_{j,a}$; $M_a$ & Task order of vehicle $j$; station of vehicle $j$ for alert $a$; vehicle pair observing alert $a$ & ---; m; --- \\
  $t^{\mathrm{rd}}_{v,k}$; $t_a^{\mathrm{start}}$ & Ready time of the $k$-th task of vehicle $v$; observation start time of alert $a$ & s \\
  $P,\hat P,\delta$; $s$; $\sigma(s)=a+bs$ & True position, estimated position, and localization error; cumulative travel distance; single-axis localization standard deviation ($a,b\ge0$) & m; m; m, m/m \\
  $\varepsilon_0,\varepsilon_c$; $\Delta\psi$; $\sigma_\psi$; $s_{\max}$ & Gimbal bias/control error; difference between commanded and true azimuth ($\wrap$); combined standard deviation of $\Delta\psi$; maximum travel distance satisfying $P(|\Delta\psi|\le5^\circ)\ge0.95$ & $0.8^\circ/0.6^\circ$; $^\circ$; $^\circ$; m \\
  $t_p,t_{\mathrm{rec}}$; $N,k,\eta$; $t_{\mathrm{last}}$ & Rescheduling trigger/communication recovery time; total alerts, successful starts, $\eta=k/N$; latest observation start time (60~s if not started) & s; ---; s \\
  \bottomrule
\end{tabularnb}
}
\end{table}

\subsection{Assumptions}

\begin{enumerate}[label=(A\arabic*)]
\item \textbf{Site and building geometry.} The site is level; the five buildings are axis-aligned cuboids standing on the ground, with the given planar outlines and heights. The coordinate system and building geometry are collected in Section~4.1. Terrain, vegetation, and other obstacles are ignored.

\item \textbf{Vehicle geometry and motion.} An observation vehicle is a disk of radius $1.5$~m; its center must lie at least $1.5$~m inside the site boundary and at least $1.5$~m from each building footprint (Minkowski expansion, Section~4.1). Vehicles travel at a constant $3$~m/s; acceleration, braking, turning, and collision avoidance are ignored. Paths are shortest obstacle-avoiding paths, and straight-line distance is only a lower bound.

\item \textbf{Gimbal kinematics.} The pitch range $[-10^\circ,80^\circ]$ is closed. Both axes rotate at $30^\circ$/s and can rotate simultaneously; the slew time is the longer axis, and start/stop transients and control delay are ignored. During travel the gimbal remains level facing east; it may turn toward the target only after the vehicle stops.

\item \textbf{Line of sight and occlusion.} The LOS is the straight segment from the camera to the target; buildings are opaque. Tangency or contact counts as occlusion; the 3-D slab test is given in Section~4.1. Atmospheric refraction and rain/fog are ignored.

\item \textbf{Targets and alerts.} Targets are stationary during deployment. An alert $(\theta,R,H)$ is a nominal planning datum: its nominal values define the target position used in geometry, feasibility, and scheduling computations, but it is neither an error-free measurement nor a localization result. No alert error model is adopted because the problem specifies none; the localization-accuracy results of Problem~2 are conditional on the two-station bearing-error model, and the ``perfectly known target position'' of Problem~3 takes the nominal datum as the pointing reference, as specified. Height is used only for elevation and occlusion tests; it does not enter the intersection solution.

\item \textbf{Observation timing and information.} Two vehicles begin a $5$~s synchronized observation after arriving and completing the slew; they then reset and travel to the next station, with all phases counted on the timeline. During travel the gimbal remains east; after arrival it performs a full-angle slew. Only when a dispatched-but-unobserved vehicle is redispatched and the new station coincides with its current position (no travel segment) does the first post-snapshot task continue from the current attitude; all later tasks start from east. Except during a communication outage, commands, feedback, and computation have no delay.

\item \textbf{Error statistics.} In Problem~2, the two stations' azimuth errors are independent $\mathcal{N}(0,(1^\circ)^2)$; this $1^\circ$ budget is the azimuth measurement error at surveyed stations, absorbs any residual docking error, and does not include the Problem~3 self-localization term. In Problem~3, the vehicle's horizontal coordinate errors are independent, identically distributed normal in $x$ and $y$; the gimbal bias $0.8^\circ$ and control error $0.6^\circ$ are independent of each other and of the localization error. All sources have zero mean.

\item \textbf{Decision information and disturbances.} At $t=0$, when the baseline plan is made, no disturbance is known in advance; rescheduling is triggered only after a disturbance occurs, using real-time information. No random events other than the specified disturbances are considered. A failed vehicle exits immediately and does not block traffic.
\end{enumerate}

% =====================================================================
% 4. Model Establishment and Solution  (4.1 -- 4.3)
% =====================================================================
\section{Model Establishment and Solution}
\label{sec:model}

\subsection{Geometric Preprocessing and Conventions}

All four problems share the same site geometry and observation model. This section collects the three conventions used repeatedly below; they are not repeated in each problem.

\noindent (a) \textbf{Coordinate system and site.} The southwest corner is the origin, the $x$-axis points east, and the $y$-axis points north; the site is $[0,100]^2$, and the dispatch center is $O=(50,50)$. A building is the axis-aligned cuboid $[x_i^0,x_i^1]\times[y_i^0,y_i^1]\times[0,h_i]$ standing on $z=0$; the camera height is $h_{\mathrm{cam}}=2$~m. The horizontal projection of an alert $(\theta,R,H)$ is $T'=(50+R\cos\theta,50+R\sin\theta)$, and its height is $H$.

\noindent (b) \textbf{Vehicle geometry and Minkowski expansion.} A vehicle is a disk of radius $R_{\mathrm{veh}}=1.5$~m centered at $p=(x,y)$. The vehicle geometry constraint is
\begin{equation*}
  \mathrm{C1}=\{p:\ 1.5\le x\le98.5,\ 1.5\le y\le98.5,\ d(p,B_i)\ge1.5,\ i=1,\dots,5\},
\end{equation*}
where
\begin{equation*}
  d(p,B)=\sqrt{[\max(x_0-x,0,x-x_1)]^2+[\max(y_0-y,0,y-y_1)]^2};
\end{equation*}
this is equivalent to expanding each building footprint by $1.5$~m (Minkowski dilation).

\noindent (c) \textbf{Slab test for LOS occlusion.} The LOS is the 3-D segment from the camera $(p,2)$ to the target $T=(T',H)$. For the $k$-th axis interval $[\mathrm{lo}_k,\mathrm{hi}_k]$, compute the parameter interval $[t_k^-,t_k^+]$ and set $t_{\mathrm{enter}}=\max_k t_k^-$, $t_{\mathrm{exit}}=\min_k t_k^+$. The segment intersects the cuboid (including tangency) if and only if $\max(t_{\mathrm{enter}},0)\le\min(t_{\mathrm{exit}},1)$. Tangency counts as occlusion, and numerical intersection tests use a tolerance of $10^{-9}$~\cite{preparata}. The LOS occlusion constraint is
\begin{equation*}
  \mathrm{C3}=\bigcap_{i=1}^{5}\{p:\ L(p)\cap B_i=\varnothing\}.
\end{equation*}
The observable region is $\Omega=\mathrm{C1}\cap\mathrm{C2}\cap\mathrm{C3}$, where the gimbal pitch constraint C2 is given by $\varepsilon(p)=\arctan(\Delta z/r(p))\in[-10^\circ,80^\circ]$; for a given target, $\Delta z=H-h_{\mathrm{cam}}$, which is equivalent to $r\ge\Delta z/\tan80^\circ$, with $\Delta z$ evaluated per problem. The gimbal timing and the rule ``return to east during travel, turn after arrival'' follow A3/A6 and are not repeated below.

% ---------------------------------------------------------------------
\subsection{Problem 1: Observable Region and Optimal Station}
\label{sec:q1}

\subsubsection{Constraint Formalization and Optimization Model}

The horizontal projection and 3-D position of T1 are
\begin{equation*}
  T'=(50+40\cos30^\circ,\ 50+40\sin30^\circ)=(84.6410,70.0000)~\mathrm{m},
\end{equation*}
\begin{equation*}
  T=(84.6410,70.0000,45)~\mathrm{m},\qquad \Delta z=45-2=43~\mathrm{m}.
\end{equation*}
From C1/C3 in Section~4.1, the gimbal pitch constraint for T1 ($\Delta z=43$~m) is equivalent to
\begin{equation*}
  \mathrm{C2}=\{p:\ r(p)\ge\rmin\},\qquad \rmin=\Delta z/\tan80^\circ=43/5.6712818=7.5821~\mathrm{m}.
\end{equation*}
The observable region is $\Omega=\mathrm{C1}\cap\mathrm{C2}\cap\mathrm{C3}$. With $Q(p)=1/r(p)$ as the observation quality, the optimal station is
\begin{equation*}
  p^*=\argmax_{p\in\Omega}Q(p)=\argmin_{p\in\Omega}r(p).
\end{equation*}
The criterion is chosen for four reasons. (i) A single-station angle error $\delta\alpha$ translates into a lateral position error $r\,\delta\alpha$ at the target, so $1/r$ has the same direction as single-station accuracy. (ii) C2 itself is $r\ge\rmin$, so $Q$ is aligned with the active constraint, and maximizing $Q$ is equivalent to pushing the pitch to its upper bound. (iii) In the two-station accuracy model of 2(1), DRMS increases monotonically with the station--target distances $d_i$ (Section~4.3.1), so a smaller $r$ is necessary for two-station accuracy. (iv) $Q$ is coordinate-independent, dimensionally simple, and strictly monotone, which facilitates tie-breaking. $Q$ is not the only admissible criterion; the alternative $Q_2=(1/r)\cos^2(\varepsilon-45^\circ)$ was also tested, and both criteria have the same optimal set.

\subsubsection{Global Lower Bound and Optimality Theorem}

\begin{lemma}
For any $p\in\Omega$, $r(p)\ge\rmin$.
\end{lemma}

\noindent\textbf{Proof.} By C2, $\varepsilon(p)\le80^\circ$. Since $\arctan$ is strictly increasing on $(0,\infty)$ and $\Delta z>0$, $r(p)\ge\Delta z/\tan80^\circ=\rmin$. \hfill$\square$

\begin{theorem}
If $\Omega\cap\{p:r(p)=\rmin\}$ is nonempty, then $r^*=\rmin$, $\varepsilon^*=80^\circ$, $Q^*=1/\rmin=0.1319$~m$^{-1}$, and the optimal set is the arc(s) of the $\rmin$ circle that lie in $\Omega$.
\end{theorem}

\noindent\textbf{Proof.} By Lemma~1, $r^*\ge\rmin$. If some $p_0$ satisfies $r(p_0)=\rmin$, then $r^*\le\rmin$, hence $r^*=\rmin$. Then $\varepsilon^*=\arctan(\Delta z/\rmin)=80^\circ$ and $Q^*=1/\rmin$. Every optimizer must satisfy $r=\rmin$, so the optimal set is the intersection of $\Omega$ with the circle, i.e., a union of arcs. \hfill$\square$

The theorem also shows that $r^*$ is independent of the building layout; the environment changes only the position and length of the optimal arc. This is the origin of the principle ``the optimal value is robust, the optimal location is sensitive.''

\subsubsection{Optimal Arc and Tie-Breaking}

The distance from $T'$ to the east face $x=79$ of B2 is only $5.6410$~m, less than $\rmin$, so the $\rmin$ circle intrudes westward into the $1.5$~m buffer of B2. Substituting points on the circle into the two inequalities of the expanded rectangle $[61.5,80.5]\times[57.5,72.5]$ gives the blocked arc $\varphi\in[160.75^\circ,236.89^\circ]$ ($76.15^\circ$); the feasible arc crosses $0^\circ$ continuously and has length $283.85^\circ$ ($78.85\%$ of the circle). The derivation is in Appendix~A.1. The distances from the other buildings to $T'$ exceed $\rmin$, and LOS on the feasible arc passes east of or above B2, so C3 does not shorten the feasible arc further (proof in Appendix~A.1); 200,000-point arc sampling agrees with the analytics. In addition, the explicit feasible point $p_c=(84.641,77.582)$ has LOS with $x\equiv84.641>79$ throughout, does not intersect any building, and satisfies C1; thus both steps---``lower bound'' and ``attainability''---are verified by explicit construction.

The optimal arc is made unique by a three-level lexicographic tie-break:
\begin{enumerate}[label=(T\arabic*)]
\item Maximum-safety-margin band:
\begin{equation*}
  m(p)=\min\bigl\{\min_i[d(p,B_i)-1.5],\ x-1.5,\ 98.5-x,\ y-1.5,\ 98.5-y\bigr\},
\end{equation*}
keep the tie band $m(p)\ge m^*-0.05$~m, where $m^*=9.6746$~m ($\varphi=56.5038^\circ$).
\item Shortest reachable within the band: choose the point with the shortest obstacle-avoiding travel time from the nearest standby vehicle.
\item Coordinate lexicographic order: if still tied, choose the smallest $(x,y)$ in lexicographic order.
\end{enumerate}
The result is $p^*=(88.8753,76.2895)$~m, $\varphi=56.05^\circ$, $r^*=7.5821$~m, $\varepsilon^*=80^\circ$, $Q^*=0.1319$~m$^{-1}$. The margin of $p^*$ itself is $9.6247$~m, entering the tie band with a $0.0001$~m surplus; the travel time to the nearest vehicle V4 is $2.78$~s, and the minimum LOS clearance is $7.41$~m. It should be emphasized that $p^*$ is ``the fastest station within the near-maximum-margin band,'' not the maximum-margin point itself.

\subsubsection{Results}

A $0.1$~m midpoint grid estimates the observable area as $5610.4$~m$^2$ (about $56.1\%$ of the site); a $0.05$~m midpoint recomputation gives $5608.5$~m$^2$ (difference $1.87$~m$^2$, $0.033\%$, converged). The value $5614.9$~m$^2$ is the $0.25$~m grid benchmark of the sensitivity scan in Figure~\ref{fig:4-3} and is used only for relative comparisons under the same convention; the node-count convention $5628.7$~m$^2$ is for reference only. The minimum feasible $r$ on the grid is $7.5828$~m (at $(92.2,69.4)$), $0.0007$~m above the closed-form bound, which is explained by grid quantization. C1, C2, and C3 remove $1685.7$, $157.9$, and $2547.7$~m$^2$, respectively, leaving $5610.4$~m$^2$ (the three sum to $4391.3$~m$^2$; the $1.7$~m$^2$ difference from $10000-5610.4$ is grid quantization). The LOS passes north of the north face of B2, with a minimum clearance of $7.41$~m (see Appendix~A.1(b)). Figure~\ref{fig:4-1} shows the observable region, Figure~\ref{fig:4-2} shows the optimal arc and the margin--azimuth curve, and Figure~\ref{fig:4-3} shows the sensitivity of $r^*$ and the area to $H$, $\varepsilon_{\max}$, the B2 height, and the vehicle radius.

\begin{figure}[htbp]
  \centering
  \includegraphics[width=\figwidth]{fig4-1_region.png}
  \caption{Observable region for T1 and building occlusion shadows (including the $\rmin$ exclusion zone)}
  \label{fig:4-1}
\end{figure}

In Figure~\ref{fig:4-1}, the black circle is the $\rmin=7.5821$~m exclusion zone; the blue shaded area outside the circle and not shadowed by buildings is $\Omega$, with color depth corresponding to $r$ from small to large. The dark-blue narrow band immediately northeast of $T1'$ is the candidate-station region with the largest $Q=1/r$.

\begin{figure}[htbp]
  \centering
  \includegraphics[width=\figwidth]{fig4-2_arc_margin.png}
  \caption{Optimal arc and margin--azimuth curve ((a) feasible and blocked arcs; (b) $m(p)-m^*$)}
  \label{fig:4-2}
\end{figure}

In Figure~\ref{fig:4-2}(a), green is the feasible arc on the $\rmin$ circle, red is the arc $[160.75^\circ,236.89^\circ]$ blocked by B2, and $p^*$ is on the northeast side at $\varphi=56.05^\circ$; in (b), the margin curve has two sharp peaks near $m^*$ at $\varphi\approx56^\circ$ and $\approx310^\circ$, and tie-breaking selects the nearest reachable station within the left peak.

\begin{figure}[htbp]
  \centering
  \includegraphics[width=\figwidth]{fig4-3_sensitivity.png}
  \caption{Sensitivity of Problem~1 ((a) $r^*$--$H$; (b) $r^*$--$\varepsilon_{\max}$; (c) area--B2 height; (d) area--vehicle radius)}
  \label{fig:4-3}
\end{figure}

Figure~\ref{fig:4-3}(a)(b) shows that $r^*$ increases linearly with target height $H$ and decreases nonlinearly as the pitch limit $\varepsilon_{\max}$ rises; (c)(d) shows that the observable area shrinks first quickly and then slowly as B2 grows taller, and decreases almost linearly with the vehicle radius. Together, the four panels show that the pitch limit is the most sensitive geometric parameter for $r^*$.

\begin{table}[htbp]
  \centering
  \caption{Main results of Problem~1}
  \label{tab:4-1}
  \begin{tabular}{@{}p{55mm}p{105mm}@{}}
    \toprule
    Item & Result \\
    \midrule
    Target horizontal projection $T'$ & $(84.6410,70.0000)$~m \\
    Closed-form lower bound $r^*$ & $7.5821$~m ($43/\tan80^\circ$) \\
    Optimal elevation $\varepsilon^*$ & $80.000^\circ$ (active constraint) \\
    Optimal quality $Q^*$ & $0.1319$~m$^{-1}$ \\
    Optimal station $p^*$ & $(88.8753,76.2895)$~m, $\varphi=56.05^\circ$ \\
    Observable area & $5610.4$~m$^2$ (about $56.1\%$) \\
    Optimal arc & feasible $283.85^\circ$, blocked $[160.75^\circ,236.89^\circ]$ \\
    Grid verification & $7.5828$~m @ $(92.2,69.4)$, $+0.0007$~m \\
    Minimum LOS clearance / nearest vehicle & $7.41$~m / V4, $2.78$~s \\
    \bottomrule
  \end{tabular}
\end{table}

\subsubsection{Algorithm and Efficiency}

This problem uses ``closed-form lower bound + arc parameterization + three-level lexicographic rule'' rather than 2-D continuous optimization over $(x,y)$. The steps are: (i) solve C2 for $\rmin=\Delta z/\tan80^\circ$; (ii) sample 200,000 points uniformly on the $r=\rmin$ circle (about $0.0018^\circ$ spacing) and mark the feasible arc by 3-D slab intersection; (iii) make the choice unique on the feasible arc by ``maximum-safety-margin band---shortest obstacle-avoiding path within the band---coordinate lexicographic order''; (iv) independently recompute the area and minimum $r$ on a $0.1$~m grid (about $10^6$ points). The reason for this algorithm is that $Q=1/r$ is strictly monotone, so the optimal value must be attained on the active constraint C2 and the closed form provides a global optimality certificate. A 2-D grid search would only give the best grid value, without an optimality proof, at a cost of $10^6$ points, whereas the closed-form part is $O(1)$. The grid is used only for independent verification: arc sampling is $O(2\times10^5)$, and the grid verification is one $O(10^6)$ pass. The closed-form $7.5821$~m and the grid $7.5828$~m differ by only $0.0007$~m, confirming that the certificate and the numerical recomputation agree.

% ---------------------------------------------------------------------
\subsection{Problem 2: Two-Station Intersection and Multi-Target Scheduling}
\label{sec:q2}

\subsubsection{Two-Station Intersection Accuracy Model}

Let the two stations be $p_i$, with station--target distances $d_i$, LOS azimuths $\alpha_i$, unit directions $u_i=(\cos\alpha_i,\sin\alpha_i)$, and unit normals $n_i=(-\sin\alpha_i,\cos\alpha_i)$. The azimuth errors $\delta\alpha_i\sim\mathcal{N}(0,\sigma^2)$ are independent. The intersection-point error satisfies the first-order propagation \cite{torrieri,foy}
\begin{equation*}
  A\,\Delta p=\delta,\qquad A=\begin{bmatrix}n_1^T\\ n_2^T\end{bmatrix},\qquad
  \delta=(d_1\delta\alpha_1,d_2\delta\alpha_2)^T,
\end{equation*}
\begin{equation*}
  \Sigma=A^{-1}DA^{-T},\qquad D=\diag(d_1^2\sigma^2,d_2^2\sigma^2).
\end{equation*}
Hence $\trace\Sigma=\sigma^2(d_1^2+d_2^2)/\sin^2\gamma$ and $\det\Sigma=\sigma^4d_1^2d_2^2/\sin^2\gamma$ (full algebra in Appendix~A.2).

Define
\begin{equation*}
  G=\frac{\sqrt{(d_1^2+d_2^2)/2}}{|\sin\gamma|},\qquad
  \mathrm{DRMS}=\sigma G=\sqrt{\trace\Sigma/2};
\end{equation*}
the 2-D radial RMSE is $\sqrt2\,\mathrm{DRMS}$, and the $95\%$ circle radius is $R_{95}=2.4477\,\mathrm{DRMS}$ (exact when isotropic).

\subsubsection{Global Lower Bound and Vehicle/Station Selection}

From $|\sin\gamma|\le1$, the QM--AM (quadratic-mean--arithmetic-mean) inequality, and $d_i\ge\rmin$ from C2,
\begin{equation*}
  G\ge\sqrt{\frac{d_1^2+d_2^2}{2}}\ge\frac{d_1+d_2}{2}\ge\rmin=\frac{35-2}{\tan80^\circ}=5.8188~\mathrm{m}.
\end{equation*}
Equality requires $\gamma=90^\circ$ and $d_1=d_2=\rmin$. For T2, the geometrically feasible arc on the $\rmin$ circle is $[59.66^\circ,244.24^\circ]$; with a $0.5$~m LOS/geometric margin, the robust feasible arc is $[70.95^\circ,238.63^\circ]$, of length $167.68^\circ>90^\circ$, so a $90^\circ$ pair exists. Under the 30~s limit, the V2-related combinations have $G=8.9678$~m, but the V2-side travel alone is $30.29$~s (V1+V2) or $34.75$~s (V2+V3, V2+V4, V2+V5), which already exceeds 30~s, so all four are infeasible (V1+V2 is the thinnest, $30.29$~s, against a straight-line lower bound of $30.00$~s); paths are computed on the visibility graph with Dijkstra~\cite{dijkstra}. In the stricter $0.75$~m robust inner domain, V3+V5 is selected by ``minimum total travel time---minimum maximum single-vehicle time---lexicographic order'':
\begin{equation*}
  p_1=(5.2787,73.3912)~\mathrm{m},\quad p_2=(11.9201,78.2501)~\mathrm{m},
\end{equation*}
\begin{equation*}
  d_1=d_2=5.8188~\mathrm{m},\quad \gamma=90.000^\circ,\quad \beta=8.2290~\mathrm{m},
\end{equation*}
with travel times $5.11/13.61$~s (sum $18.72$~s), $\mathrm{DRMS}=0.1016$~m, and $R_{95}=0.2486$~m.

\begin{table}[htbp]
  \centering
  \caption{Two-station intersection accuracy and station-selection results}
  \label{tab:4-2}
  \begin{tabular}{@{}p{55mm}p{105mm}@{}}
    \toprule
    Item & Result \\
    \midrule
    Selected vehicles & V3+V5 \\
    Stations $p_1$ / $p_2$ & $(5.2787,73.3912)$ / $(11.9201,78.2501)$~m \\
    Station--target distances $d_1=d_2$ & $5.8188$~m \\
    Intersection angle $\gamma$ / baseline $\beta$ & $90.000^\circ$ / $8.2290$~m \\
    Travel times & $5.11$ / $13.61$~s (sum $18.72$~s) \\
    DRMS / $R_{95}$ & $0.1016$ / $0.2486$~m \\
    MC radial RMSE / P95 & $0.1436$ / $0.2486$~m \\
    Robust fallback ($78^\circ$) & $r\ge7.0144$~m, $\mathrm{DRMS}=0.1224$~m ($+20.5\%$) \\
    \bottomrule
  \end{tabular}
\end{table}

\begin{figure}[htbp]
  \centering
  \includegraphics[width=\figwidth]{fig4-4_t2_region_reach.png}
  \caption{Observable region of T2 and the 30~s reachable domains ((a) full domain; (b) zoom at the selected stations)}
  \label{fig:4-4}
\end{figure}

Figure~\ref{fig:4-4}(a) shows the observable region of T2 and the 30~s reachable domains of the five vehicles (dashed circles of radius 90~m). The target $T2'$ lies above B3, and the feasible stations are mainly on the west side of the building. The zoom in (b) shows the selected stations $p_1=(5.2787,73.3912)$~m and $p_2=(11.9201,78.2501)$~m; both are $5.8188$~m from the target, the intersection angle is $\gamma=90.000^\circ$, and the baseline is $8.2290$~m. The travel times are $5.11$~s and $13.61$~s, both within the 30~s limit.

\begin{figure}[htbp]
  \centering
  \includegraphics[width=\figwidth]{fig4-5_accuracy_params.png}
  \caption{Two-station accuracy versus parameters ((a) DRMS--$\gamma$; (b) DRMS--$d$; (c) DRMS--baseline)}
  \label{fig:4-5}
\end{figure}

Figure~\ref{fig:4-5}(a) shows that DRMS is U-shaped in the intersection angle $\gamma$, minimized at $\gamma=90^\circ$ and diverging rapidly as $\gamma\to0^\circ/180^\circ$ (the shaded region is degenerate); the selected point ($\gamma=90^\circ$, $d=5.8$~m) has $\mathrm{DRMS}=0.1016$~m. (b) shows that DRMS grows almost linearly with the station--target distance $d$, with the smallest slope at $\gamma=90^\circ$; when $d$ increases from $5.8$~m to $30$~m, DRMS rises to about $0.75$~m. (c) shows that, at a fixed intersection angle, DRMS increases linearly with the baseline $B$, with the steepest slope at $\gamma=60^\circ$.

\begin{figure}[htbp]
  \centering
  \includegraphics[width=\figwidth]{fig4-6_error_ellipses.png}
  \caption{Intersection error ellipses and 95\% circle (with a local zoom)}
  \label{fig:4-6}
\end{figure}

The zoom in Figure~\ref{fig:4-6} shows the geometry of the selected stations $p_1$, $p_2$, and $T2'$; the local inset gives the $1\sigma$, $2\sigma$, and $95\%$ position-error circles at the target. Because the two stations are equidistant and orthogonal, the error distribution is nearly isotropic with $R_{95}=0.2486$~m. The error-circle radii are much smaller than the station--target distances, indicating that the intersection geometry of this station pair is close to ideal.

Degeneracy analysis: as $\gamma\to0^\circ/180^\circ$, $G\to\infty$ and the condition number is $\kappa=1/|\sin\gamma|$; at $\gamma=30^\circ$, $G=2G^*$, so $[30^\circ,150^\circ]$ is used as the operating window, and a ridge-regularized least-squares estimator is used on the estimation side. The optimal station lies on the C2 boundary ($\varepsilon=80^\circ$); if a pitch margin is required, one may fall back to $\varepsilon\le78^\circ$, giving $r\ge7.0144$~m and $\mathrm{DRMS}=0.1224$~m ($+20.5\%$).

\subsubsection{Six-Slot, Five-Vehicle Scheduling Model}

Let the alert set be $A=\{\mathrm{T1},\mathrm{T2},\mathrm{T3}\}$ and the vehicle set be $V=\{\mathrm{V1},\dots,\mathrm{V5}\}$. Each alert $a$ is assigned two distinct vehicles $M_a$; the task sequence of vehicle $v$ is $\pi_v$, and the station of task $k$ is $p_{v,k}\in\Omega_a$. For task $k$, let $D_{v,k-1}$ be the previous task's departure time and $\mathrm{pos}_{v,k-1}$ the previous position. Then
\begin{equation*}
  \tau_{\mathrm{tr}}=\distpath(\mathrm{pos},p)/3,\qquad
  \tau_{\mathrm{sl}}=\max(|\alpha|,|\varepsilon|)/30,
\end{equation*}
\begin{equation*}
  t^{\mathrm{rd}}_{v,k}=D_{v,k-1}+\tau_{\mathrm{tr}}+\tau_{\mathrm{sl}},
\end{equation*}
\begin{equation*}
  S_a=\max(t^{\mathrm{rd}}_{v,k},t^{\mathrm{rd}}_{v',k'}) \quad\text{(the later of the two vehicles' ready times)},
\end{equation*}
\begin{equation*}
  D_{v,k}=S_a+5+\tau_{\mathrm{rs}},\qquad
  \tau_{\mathrm{rs}}=\max(|\alpha|,|\varepsilon|)/30.
\end{equation*}
The objective is $\min\max_a S_a$. Each alert has 10 vehicle pairs, so the three alerts give $10^3=1000$ assignment structures. For each structure, reuse orders are enumerated; stations are initialized greedily, refined by Gauss--Seidel coordinate descent and Nelder--Mead~\cite{nelder}; if two reused vehicles wait for each other and the fixed-point iteration makes no progress, the structure is declared deadlocked and removed.

Lower bound: $\mathrm{LB}(v,a)=\min_p[\distpath(V_v^0,p)/3+\max(|\alpha|,|\varepsilon|)/30]$, $S_a\ge$ the second-smallest $\mathrm{LB}(v,a)$, and the global $\mathrm{LB}=\max_a$ second-smallest $=3.414$~s. This bound ignores reuse and synchronization, so it is loose.

\subsubsection{Nearest-Assignment Baseline and Comparison}

In the order T1$\to$T2$\to$T3, fill two slots for each alert. For each slot, among vehicles not already serving this alert, choose the one with the smallest ``travel time from the current chain end to its minimum-ready station for this alert'' (ties: smaller vehicle number), and use that minimum-ready station; if the vehicle already has tasks, append to its chain. The same timeline simulation computes the start times. This baseline considers only ``who is nearest'' and ignores synchronized pairing, reuse-chain order, and slew angle.

Optimized plan: T1$\leftarrow$(V3,V5), T2$\leftarrow$(V1,V4), T3$\leftarrow$(V2,V3), with the V3 reuse chain T1$\to$T3; the start times are $1.607/6.994/10.355$~s and the objective is $10.355$~s. Baseline: T1$\leftarrow$(V3,V5), T2$\leftarrow$(V1,V3), T3$\leftarrow$(V1,V5), with reuse chains for V1, V3, and V5; the start times are $1.607/11.140/20.968$~s and the objective is $20.968$~s. The latest start is reduced by $10.613$~s ($-50.6\%$).

The difference has three sources. (i) Synchronized-pairing loss: the baseline places T2 as the second task in V3's reuse chain, so V3 must first finish T1 and reset, forcing V1 on the other side of T2 to wait; the optimized plan instead pairs V1 with the dedicated V4. (ii) Reuse-chain order: the baseline's three two-task chains serialize slow tasks after fast alerts and block T3's early window; the optimized plan has only one T1$\to$T3 chain. (iii) Azimuth-driven station selection: ready time is highly sensitive to the pan angle ($30^\circ$/s), while the baseline selects stations by distance and ignores the turn angle; the optimization model minimizes ready time.

Note that the T2 pair of 2(1) is accuracy-optimal, whereas 2(2) re-optimizes the stations for timeliness; Section~4.3.6 verifies the accuracy that this costs.

\subsubsection{Results and Robust Deployment}

The $10.355$~s optimum has knife-edge station clearances of only $1.3\times10^{-3}$~m (V3 at T3) and $6.7\times10^{-2}$~m (V4 at T2); the full stations and timings are listed in Table~\ref{tab:D-2}. Re-optimizing with $0.5$~m clearance yields a robust variant of $10.837$~s ($+4.7\%$), with assignment T1$\leftarrow$(V3,V5), T2$\leftarrow$(V1,V3), T3$\leftarrow$(V4,V5) and start times $1.54/10.804/10.837$~s; the robust variant is recommended for deployment. The sensitivities to speed and angular rate are given in Section~6.2; the elasticity of $\omega$ is about four times that of $v$.

\begin{table}[htbp]
  \centering
  \caption{Scheduling plan versus the nearest-assignment baseline}
  \label{tab:4-3}
  \begin{tabular}{@{}p{38mm}p{55mm}p{65mm}@{}}
    \toprule
    Item & Optimized plan & Nearest-assignment baseline \\
    \midrule
    T1 start & $1.607$~s & $1.607$~s \\
    T2 start & $6.994$~s & $11.140$~s \\
    T3 start & $10.355$~s & $20.968$~s \\
    Objective max & $10.355$~s & $20.968$~s \\
    Assignment & T1(V3,V5), T2(V1,V4), T3(V2,V3) & T1(V3,V5), T2(V1,V3), T3(V1,V5) \\
    Reuse chains & V3: T1$\to$T3 & V1: T2$\to$T3; V3: T1$\to$T2; V5: T1$\to$T3 \\
    \bottomrule
  \end{tabular}
\end{table}

The baseline Gantt chart is structurally the same as the optimized plan and is omitted (its times are listed in Table~\ref{tab:4-3}).

\begin{figure}[htbp]
  \centering
  \includegraphics[width=\figwidth]{fig4-7_gantt_inset.png}
  \caption{Vehicle Gantt chart and deployment routes of the optimized plan ((a) Gantt chart; (b) routes and stations)}
  \label{fig:4-7}
\end{figure}

Figure~\ref{fig:4-7}(a) gives the complete task chain of the optimized plan on the time axis: T1, T2, and T3 begin synchronized two-vehicle observation at $1.607$, $6.994$, and $10.355$~s, respectively; V3 completes T1 and, after reset and slew, continues to T3, forming the only reuse chain. The orange slew segments and gray synchronized-waiting segments explain the composition of each alert's start time. (b) shows the post-arrival deployment positions of the five vehicles and the three alert bearings: T1 is observed by V3 and V5 from the same northwest side ($\gamma$ only $19.34^\circ$), T2 is bracketed by V1 and V4 from the south and east ($\gamma=85.08^\circ$), and T3 is bracketed by V2 and V3 from the southeast and northwest ($\gamma=127.39^\circ$), providing the geometric basis for the accuracy check in Section~4.3.6.

\subsubsection{Accuracy Check of the Timeliness Plan}

The stations of 2(2) are re-optimized independently for timeliness and must be substituted back into the accuracy model of 2(1) to check the intersection geometry. Using $G$ and DRMS from Section~4.3.1, the three station pairs are evaluated in Table~\ref{tab:4-4}. T2 and T3 have $\gamma=85.08^\circ$ and $127.39^\circ$, inside the $[30^\circ,150^\circ]$ operating window; T1 has $\gamma=19.34^\circ$ ($\kappa=3.02$) and $\mathrm{DRMS}=3.084$~m, approaching degeneracy---the two stations are on the same side of the target and the LOS are nearly parallel. If T1 must also satisfy $\gamma\ge30^\circ$, fixing the assignment and order and reselecting T1's two stations gives $(14.90,85.60)/(53.90,98.40)$, $\gamma=30.12^\circ$, $G=116.68$~m, $\mathrm{DRMS}=2.036$~m; T1's start is delayed from $1.607$~s to $3.746$~s, and the maximum start over the three alerts rises from $10.355$~s to $12.983$~s ($+2.629$~s, $+25.38\%$). This shows that timeliness and accuracy are not interchangeable: if deployment requires T1 to satisfy $\gamma\ge30^\circ$, an intersection-angle constraint must be added to the scheduling model.

\begin{table}[htbp]
  \centering
  \caption{Intersection geometry and accuracy of the timeliness plan}
  \label{tab:4-4}
  \begin{tabular}{@{}p{18mm}p{25mm}p{25mm}p{35mm}p{35mm}@{}}
    \toprule
    Alert & Pair & $\gamma$ ($^\circ$) & $G$ (m) & DRMS (m) \\
    \midrule
    T1 & V3+V5 & 19.34 & 176.67 & 3.084 \\
    T2 & V1+V4 & 85.08 & 67.93 & 1.186 \\
    T3 & V2+V3 & 127.39 & 81.32 & 1.419 \\
    T1* & V3+V5 & 30.12 & 116.68 & 2.036 \\
    \bottomrule
  \end{tabular}

  \vspace{2pt}
  {\footnotesize Note: $\mathrm{DRMS}=\sigma G$, $\sigma=1^\circ=0.017453$~rad; T1* is the constrained re-solution with $\gamma_{\mathrm{T1}}\ge30^\circ$ under the same assignment/order and does not constitute a global optimality claim. None of the four station pairs is isotropic (the eigenvalue ratio of T1 is about $46:1$); exact $95\%$ coverage should use the ellipse semi-axes $\sqrt{5.991\lambda_i}$, which are $(10.56,1.55)/(3.27,2.48)/(4.77,1.19)/(6.87,1.58)$~m for T1/T2/T3/T1*, respectively.}
\end{table}

\subsubsection{Algorithm and Efficiency}

2(1) uses ``covariance propagation + analytical lower-bound certificate + pair enumeration.'' First, $G$ is derived from $\Sigma=A^{-1}DA^{-T}$; the QM--AM inequality and C2 then prove $G\ge\rmin$, and the equality condition is reduced to a one-dimensional test: whether a feasible pair with a $90^\circ$ intersection angle exists on the $\rmin$ circle. Then candidate stations are enumerated and refined by Nelder--Mead for all $C(5,2)=10$ vehicle pairs (4-D NM, $\le300$ iterations per call). The analytical certificate reduces the global optimality check to a one-dimensional feasibility test, avoiding local optima while keeping the computation at the scale of 10 pairs; a direct 4-D search has neither an optimality proof nor protection against local optima.

2(2) uses ``complete structural enumeration + fixed-point timeline simulation + coordinate descent/Nelder--Mead.'' The three alerts each have $C(5,2)=10$ vehicle pairs, giving $10^3=1000$ assignment structures; each structure enumerates reuse-chain orders, for a total of $5{,}880$ (structure, order) combinations, each simulated deterministically and checked for deadlock. Stations are initialized greedily, refined by Gauss--Seidel coordinate descent and Nelder--Mead, and the top 12 structures are re-refined. With only 5 vehicles and 6 slots, all $5{,}880$ deterministic simulations complete in one run, exhaust the structures, and guarantee reproducibility; a randomized search such as a genetic algorithm would neither guarantee the same solution nor explain why the plan is optimal. In terms of efficiency, $10.355$~s is $50.6\%$ lower than the nearest-assignment baseline of $20.968$~s, at a cost of only a few thousand simulations, at most 2 tasks per vehicle, and at most 160 NM iterations per call.

% =====================================================================
% 4.4 Problem 3 and 4.5 Problem 4
% =====================================================================
\subsection{Problem 3: Localization and Pointing Error Propagation}
\label{sec:q3}

\subsubsection{Error Propagation Model}

Take the true position $P$ as the origin and the true LOS as the $x'$ axis, with the target at $(d,0)$, $d=40$~m. Decompose the localization error into cross-range $\delta_\perp$ and along-range $\delta_\parallel$; under isotropy, $\delta_\perp\sim\mathcal{N}(0,\sigma(s)^2)$. The exact deviation from the estimated position to the target is $\Delta\alpha=\operatorname{atan2}(-\delta_\perp,d-\delta_\parallel)\approx-\delta_\perp/d$ (residual analysis in Appendix~A.3). Adding the gimbal bias $\varepsilon_0\sim\mathcal{N}(0,0.8^2)$ and control error $\varepsilon_c\sim\mathcal{N}(0,0.6^2)$,
\begin{equation}
  \Delta\psi=\wrap\bigl(\Delta\alpha\cdot180/\pi+\varepsilon_0+\varepsilon_c\bigr),
\end{equation}
where $\wrap(\cdot)$ folds angles into $(-180^\circ,180^\circ]$. The three terms are independent zero-mean normals, and $\wrap$ is not triggered at this problem's scale, so
\begin{equation}
  \Delta\psi\sim\mathcal{N}\bigl(0,\sigma_\psi^2(s)\bigr),\qquad
  \sigma_\psi^2(s)=\left(\frac{\sigma(s)}{40}\cdot\frac{180}{\pi}\right)^2+0.8^2+0.6^2
  =(1.432394\sigma(s))^2+1.0.
\end{equation}

\subsubsection{$s_{\max}$ and $(a,b)$ Constraints}

From $P(|\Delta\psi|\le5^\circ)=2\Phi(5/\sigma_\psi)-1\ge0.95$, we obtain $\sigma_\psi\le5/z_{0.975}=2.551067^\circ$. Substituting into the expression above,
\begin{equation}
  \sigma(s)\le\frac{\sqrt{2.551067^2-1.0}}{1.432394}=1.638446~\mathrm{m},
\end{equation}
\begin{equation}
  s_{\max}=\frac{1.638446-0.3}{0.02}=66.9223~\mathrm{m}\approx66.92~\mathrm{m}.
\end{equation}
At $s=80$~m, $\sigma(80)=1.9$~m, $\sigma_\psi=2.89945^\circ$, and $P=0.91538<0.95$. If the criterion must still hold at 80~m, then $a+80b\le1.638446$~m, which together with $a,b\ge0$ forms a triangular feasible region ($b=0$ gives $a\le1.6384$~m; $a=0$ gives $b\le0.0204806$~m/m). The nominal point $(0.3,0.02)$ lies outside; one must reduce $b$ to $0.0167306$~m/m ($-16.3\%$) or reduce $a$ to $0.0384$~m.

\subsubsection{Contribution Decomposition and Exact Verification}

Let $\mathrm{var}_{\mathrm{loc}}=(1.432394\sigma(s))^2$, $\mathrm{var}_{\mathrm{bias}}=0.64$, and $\mathrm{var}_{\mathrm{ctrl}}=0.36$. At $s=0$ the three shares are $15.6\%\bsl54.0\%\bsl30.4\%$; at $s=40$ they are $71.3\%\bsl18.4\%\bsl10.3\%$; at $s=80$ they are $88.1\%\bsl7.6\%\bsl4.3\%$; the bias:control ratio is always $64:36$. The dominant-factor switch point is obtained from $\mathrm{var}_{\mathrm{loc}}=\mathrm{var}_{\mathrm{bias}}+\mathrm{var}_{\mathrm{ctrl}}$ as $s=19.907$~m (shares $50.0\%\bsl32.0\%\bsl18.0\%$), after which localization becomes the first source. The linearization residual RMS is $\sigma(s)^2/d^2$, which at $s=80$~m is $0.1293^\circ$, only $4.46\%$ of $\sigma_\psi$ and about $0.2\%$ of the variance. A $256\times256$ Hermite quadrature gives the exact $s_{\max}=66.84$~m, $P(66.92)=0.94980$, and $a+80b\le1.6367$~m; the recommended practice is to use $66.84$~m as the conservative design value and $66.92$~m as the closed-form primary answer. The $P$ deviation of $10^6$ Monte Carlo runs \cite{hammersley} is $\le3\times10^{-4}$, validating the linearization and the normal composition.

\noindent\textbf{Deployment implications.} (i) The travel budget should be capped near the conservative $66.84$~m rather than $66.92$~m, because the closed form slightly understates the risk at the boundary. (ii) If the sensing budget is limited, reducing the slope $b$ of the localization error is far more effective than reducing the offset $a$: at 80~m one unit of $b$ saves 80 units of $\sigma$, whereas $a$ enters one-to-one. (iii) Improving the gimbal (bias plus control, $1.0^\circ$ combined) helps mainly at short travel; beyond about 100~m the localization term dominates and gimbal upgrades saturate (Section~6.2).

\begin{table}[htbp]
  \centering
  \caption{Monte Carlo verification for Problem~3 ($10^6$ samples)}
  \label{tab:4-5}
  \begin{tabular}{@{}p{16mm}p{22mm}p{25mm}p{25mm}p{30mm}p{30mm}@{}}
    \toprule
    $s$ (m) & $\sigma(s)$ (m) & $P$ empirical & $P$ analytical & $\sigma$ empirical ($^\circ$) & $\sigma$ analytical ($^\circ$) \\
    \midrule
    0 & 0.30 & 0.999996 & 0.999996 & 1.0875 & 1.0884 \\
    20 & 0.70 & 0.999584 & 0.999586 & 1.4163 & 1.4161 \\
    40 & 1.10 & 0.992694 & 0.992622 & 1.8651 & 1.8662 \\
    66.92 & 1.6384 & 0.949730 & 0.950000 & 2.5528 & 2.5511 \\
    80 & 1.90 & 0.915229 & 0.915376 & 2.9006 & 2.8995 \\
    \bottomrule
  \end{tabular}
\end{table}

\begin{figure}[htbp]
  \centering
  \includegraphics[width=\figwidth]{fig4-8_variance.png}
  \caption{Variance components and shares of $\sigma_\psi$ ((a) component curves; (b) stacked shares)}
  \label{fig:4-8}
\end{figure}

In Figure~\ref{fig:4-8}(a), the localization term rises linearly with $s$, and its intersection with the allowable line $2.551^\circ$ is $s_{\max}=66.92$~m; (b) shows that after $s\approx19.9$~m the localization share exceeds the sum of the gimbal bias and control shares, reaching $88.1\%$ at $s=80$~m.

\begin{figure}[htbp]
  \centering
  \includegraphics[width=\figwidth]{fig4-9_probability_ab.png}
  \caption{$P(|\Delta\psi|\le5^\circ)$ curve and the $(a,b)$ feasible region ((a) $P$--$s$; (b) $(a,b)$ region)}
  \label{fig:4-9}
\end{figure}

Figure~\ref{fig:4-9}(a) gives the decreasing curve of $P(|\Delta\psi|\le5^\circ)$ versus $s$; its intersection with the $0.95$ line is $s_{\max}=66.92$~m (exact $66.84$~m), and the Monte Carlo points coincide with the analytical curve. In (b), the shaded area is the $(a,b)$ feasible region where $P\ge0.95$ at 80~m; the nominal point $(0.3,0.02)$ lies outside.

\subsubsection{Algorithm and Efficiency}

This problem uses ``closed-form first-order propagation + monotone criterion + exact quadrature verification.'' The steps are: (i) linearize the lateral projection to obtain $\Delta\psi\sim\mathcal{N}(0,(1.432394\sigma(s))^2+1.0)$; (ii) from the two-sided criterion $P(|\Delta\psi|\le5^\circ)\ge0.95$, obtain $\sigma_\psi\le2.5511^\circ$ and invert for $\sigma(s)$ and $s_{\max}$; (iii) evaluate the exact $\operatorname{atan2}$ geometry by a $256\times256$ Gauss--Hermite quadrature over the two-dimensional expectation and find $s_{\max}$ by bisection; (iv) independently verify with $10^6$ Monte Carlo runs. The closed form is chosen because $\sigma_\psi(s)$ is monotone in $s$, so the criterion can be inverted directly and gives the primary answer in $O(1)$; Hermite quadrature is chosen because 65,536 nodes already reach $10^{-3}$ accuracy, at about $1/15$ of the cost of $10^6$ MC runs and without sampling error. The results are $66.92$~m in closed form and $66.84$~m exact ($0.13\%$ difference), with an MC $P$ deviation no larger than $3\times10^{-4}$.

% ---------------------------------------------------------------------
\subsection{Problem 4: Rescheduling under Disturbances}
\label{sec:q4}

\subsubsection{State Snapshot and Locking Rules}

At $t_p$ (5~s without an outage, $t_{\mathrm{rec}}$ with an outage), record $(t_{\mathrm{avail}},\mathrm{pos},\mathrm{att},\text{\texttt{locked\_until}})$ for each vehicle. Traveling: position is interpolated linearly, attitude is east, and the vehicle can be redispatched immediately. Slewing: the vehicle is at its station, attitude is interpolated, and it can be redispatched. Synchronized waiting: the vehicle is at its station, attitude points at the target, and it can be redispatched. Observing/resetting: the alert is locked and counted as served---it is not rescheduled even though its 5~s observation may still be in progress---and the vehicle remains locked until its reset finishes.

Rules: alerts that started before $t_p$ cannot be rescheduled; dispatched-but-unobserved vehicles can be redispatched; a failed vehicle is removed and its pending alerts become orphans; no new commands are received during an outage, and rescheduling is triggered after $t_{\mathrm{rec}}$. \textbf{``Continue-slew'' applies only to the first task of an unobserved redispatch}: only if the vehicle has not started observing and the new station coincides with its current position (distance $\le5$~cm), the first post-snapshot task continues from the current attitude, with
\begin{equation}
  \mathrm{slew}=\max\bigl(|\wrap(\alpha_{\mathrm{new}}-\alpha_{\mathrm{cur}})|,\ |\varepsilon_{\mathrm{new}}-\varepsilon_{\mathrm{cur}}|\bigr)/30;
\end{equation}
after that task resets, all subsequent tasks perform a full-angle slew from east. Other redispatches (new station more than 5~cm from the current position) follow A3/A6: return to east during travel and perform a full slew after arrival. For example, in the $(\theta=0^\circ,\rho=30~\mathrm{m})$ scenario (internal id \texttt{th0\_R30}), V1's redispatch chain is T2$\to$T3: T2 has a $0$~s continue-slew; after T2 observation and reset, T3 slews from east for $1.655$~s, so T3 starts at $16.772$~s. Continue-slew is the continuation of the same pointing phase, not a new ``arrival--pointing'' process.

\subsubsection{Rescheduling Model and Benchmark Strategy}

The pending alert set $A_p$ is the union of alerts whose baseline start is $\ge t_p$, orphan alerts, and the new alert T4; the available vehicle set $V_p$ contains all non-failed vehicles. The decisions are $M_a$, $\pi_v$, and $p_{v,k}$; the objective is lexicographic: maximize $k$, then minimize $t_{\mathrm{last}}$; the constraints are two distinct available vehicles per pending alert, stations in the corresponding $\Omega_a$, timeline propagation from the snapshot state, and exclusion of locked alerts and vehicles. The solution reuses the structural enumeration, coordinate descent, Nelder--Mead, and deadlock detection of Problem~2(2), changing only the initial conditions.

Keep-original: no new commands are issued. If a failed vehicle belongs to the pair of a pending alert, that alert never starts, so $k$ decreases by 1 and $t_{\mathrm{last}}=60$; the new alert has no assignment, so $N$ increases by 1 while $k$ is unchanged; with an outage only, the two strategies are equivalent.

\subsubsection{Scenario Set and Results}

The scenario set has 74 scenarios: A, new alert, $8$ azimuths $\times$ $3$ distances $=24$; B, single vehicle failure, 5; C, single outage, 3; D, failure $\times$ outage, 15; E, new alert $\times$ failure, 6; F, new alert $\times$ outage, 9; G, triple combination, 12. New alerts use 8 azimuths at $45^\circ$ and $\rho\in\{30,40,50\}$~m, covering the shadow directions of all buildings and both ends of the distance range; a subsequent $22.5^\circ$ refinement moves $t_{\mathrm{last}}$ smoothly within $14.57$--$16.45$~s, which confirms that $45^\circ$ sampling suffices. The failure class runs all 5 vehicles; the combination class uses V2 and V4 as representatives, because V1 and V4 are the two members of T2's pair and V2, V3 are the two members of T3's pair, so a failure within a pair leaves an isomorphic orphan structure (V3's only other task, T1, is locked at $t_p$), while V5's only task T1 is already locked and loses nothing. Failure $\times$ outage covers all three $t_{\mathrm{rec}}$, while the triple class uses only $t_{\mathrm{rec}}=10/30$~s: $t_{\mathrm{last}}$ responds in two stages (chain-limited, then unit slope), the middle value 20~s is already covered by D, and the secant slope ($\approx0.6$) makes the missing points predictable. The set covers azimuth, distance, failure role, outage length and combination, and its representativeness is supported by the $22.5^\circ$ refinement and the $0.5$~m margin/tolerance/robust-baseline checks.

Results: rescheduling strictly improves $\eta$ in 67 of 74 scenarios, ties in 7, and degrades none (reported at $0.01$~s precision); Table~\ref{tab:4-6} summarizes the outcome by scenario class and Table~\ref{tab:4-7} lists three representative scenarios. Representative values: V4 failure, keep $(0.667,60)\to$ resched $(1.000,11.17)$; new alert $\theta=0^\circ$, $\rho=30$~m, keep $(0.750,60)\to(1.000,16.77)$; triple combination (median new alert + V4 failure + $t_{\mathrm{rec}}=10$~s), keep $(0.500,60)\to(1.000,25.01)$. By failed vehicle, the rescheduled $t_{\mathrm{last}}$ for V1/V4, V2/V3, and V5 is $11.17$, $10.90$, and $10.35$~s, respectively, all with $\eta=1$. For a V1 failure, failure $\times$ outage with $t_{\mathrm{rec}}=10\bsl20\bsl30$ gives $t_{\mathrm{last}}=21.01\bsl25.08\bsl35.08$~s (increments $4.07\bsl10.00$~s), showing the two-stage behavior. The mechanism is ``structural loss $\to$ temporal loss''; its upper limit is jointly constrained by the azimuth of idle vehicles, the slew + travel recovery delay, and the outage duration.

\begin{table}[htbp]
  \centering
  \caption{Summary by scenario class}
  \label{tab:4-6}
  \begin{tabular}{@{}lccc@{}}
    \toprule
    Scenario class & Scenarios & $\eta$ improved & Tied \\
    \midrule
    A New alert & 24 & 24 & 0 \\
    B Single vehicle failure & 5 & 4 & 1 \\
    C Single outage & 3 & 0 & 3 \\
    D Failure $\times$ outage & 15 & 12 & 3 \\
    E New alert $\times$ failure & 6 & 6 & 0 \\
    F New alert $\times$ outage & 9 & 9 & 0 \\
    G Triple combination & 12 & 12 & 0 \\
    Total & 74 & 67 & 7 \\
    \bottomrule
  \end{tabular}

  \vspace{2pt}
  {\footnotesize Note: G uses 3 representatives $\times$ 2 failed vehicles $\times$ 2 recovery times ($t_{\mathrm{rec}}=10/30$~s only); the intermediate value $t_{\mathrm{rec}}=20$~s is covered by D.}
\end{table}

\begin{table}[htbp]
  \centering
  \caption{Comparison of representative scenarios}
  \label{tab:4-7}
  \begin{tabular}{@{}p{55mm}p{50mm}p{50mm}@{}}
    \toprule
    Scenario & Keep original $(\eta,t_{\mathrm{last}})$ & Reschedule $(\eta,t_{\mathrm{last}})$ \\
    \midrule
    V4 failure & $(0.667,60)$ & $(1.000,11.17)$ \\
    New alert $\theta=0^\circ$, $\rho=30$~m & $(0.750,60)$ & $(1.000,16.77)$ \\
    Triple (median + V4 + $t_{\mathrm{rec}}=10$~s) & $(0.500,60)$ & $(1.000,25.01)$ \\
    \bottomrule
  \end{tabular}
\end{table}

\begin{figure}[htbp]
  \centering
  \includegraphics[width=\figwidth]{fig4-10_tlast_response.png}
  \caption{Two-stage response of $t_{\mathrm{last}}$ to $t_{\mathrm{rec}}$ ((a) $t_{\mathrm{last}}$--$\theta$; (b) $t_{\mathrm{last}}$--$t_{\mathrm{rec}}$)}
  \label{fig:4-10}
\end{figure}

Figure~\ref{fig:4-10}(a) shows that the new alert's azimuth has little effect on $t_{\mathrm{last}}$ (about $15$--$17$~s), far below the 60~s of keeping the original plan; (b) shows that $t_{\mathrm{last}}$ rises piecewise with the recovery time $t_{\mathrm{rec}}$, that V5 stays at $10.36$~s because T1 is already locked and completed, and that the other failed vehicles exhibit the ``chain-limited---unit-slope'' two-stage behavior.

\begin{figure}[htbp]
  \centering
  \includegraphics[width=\figwidth]{fig4-11_resched_compare.png}
  \caption{Rescheduling Gantt chart and 74-scenario comparison ((a) V4 failure Gantt chart; (b) keep/resched scatter)}
  \label{fig:4-11}
\end{figure}

Figure~\ref{fig:4-11}(a) shows that after V4 fails at $t=5$~s, rescheduling still starts T2 and T3 at about $7$~s and $10$~s and keeps $\eta=1.00$; in (b), among the 74 scenarios, all green hollow points lie at $\eta=1.00$, the red crosses are distributed between $0.50$ and $1.00$, and the arrows show that rescheduling trades a larger $t_{\mathrm{last}}$ for a restored completion rate.

\subsubsection{Algorithm and Efficiency}

This problem uses ``event-triggered snapshot + reuse of the 2(2) solver + deterministic scenario enumeration.'' The steps are: (i) at $t_p$, record $(t_{\mathrm{avail}},\mathrm{pos},\mathrm{att},\text{\texttt{locked\_until}})$ for each vehicle and classify states as redispatchable or locked; (ii) using the snapshot as the initial condition, re-optimize the pending alerts with the structural enumeration, fixed-point simulation, coordinate descent, and Nelder--Mead of 2(2); (iii) compute $(\eta,t_{\mathrm{last}})$ for keep-original and reschedule in each of the 74 representative scenarios. This approach is chosen because the post-disturbance scale is the same order as 2(2) and reusing the same solver keeps the conventions consistent; with only 74 scenarios, deterministic enumeration completes in one run without online stochastic optimization. Keep-original is evaluated analytically and serves as a comparison with per-scenario simulation. The result is a strict $\eta$ improvement in 67/74 scenarios, 7 ties, and 0 degradations.

% =====================================================================
% 5. Model Validation and Result Analysis
% =====================================================================
\section{Model Validation and Result Analysis}
\label{sec:validation}

\begin{table}[htbp]
  \centering
  \caption{Summary of validation for the four problems}
  \label{tab:5-1}
  \small
  \begin{tabular}{@{}p{28mm}p{52mm}p{55mm}p{17mm}@{}}
    \toprule
    Object & Method & Deviation & Details \\
    \midrule
    Q1 global optimum & Closed-form $r^*$ vs.\ $0.1$~m grid & $7.5821$ vs.\ $7.5828$~m ($+0.0007$~m, $0.009\%$) & 4.2.4 \\
    Q2 accuracy and scheduling & Covariance vs.\ 200,000 MC; structural enumeration + top-5 + $0.5$~m variant & RMSE $0.1436$~m $=\sqrt2\,\mathrm{DRMS}$; P95 $0.2486$~m; optimum $10.355$~s, robust $10.837$~s & 4.3.2/4.3.5 \\
    Q3 error propagation & Linearization vs.\ $256\times256$ Hermite; $10^6$ MC & $s_{\max}$ $66.92$ vs.\ $66.84$~m ($0.13\%$); $P$ deviation $\le3\times10^{-4}$ & 4.4.3 \\
    Q4 rescheduling & 74-scenario accounting audit; $0.5$~m margin/tolerance/$22.5^\circ$ densification & 0 anomalies; $\eta/t_{\mathrm{last}}$ conclusions unchanged & 4.5.3 \\
    \bottomrule
  \end{tabular}
\end{table}

\subsection{Validation of Problem 1}

The closed-form lower bound $r^*=7.5821$~m differs from the $0.1$~m grid minimum $7.5828$~m by only $0.0007$~m ($0.009\%$); the observable area converges from $5610.4$ to $5608.5$~m$^2$ between the $0.1$~m and $0.05$~m midpoint grids ($0.03\%$ apart).

\subsection{Validation of Problem 2}

The covariance model gives $\mathrm{DRMS}=0.1016$~m; 200,000 MC runs give a radial RMSE of $0.1436$~m $=\sqrt2\,\mathrm{DRMS}$ and P95 $=0.2486$~m. The scheduling optimum $10.355$~s and the robust variant $10.837$~s corroborate each other.

\subsection{Validation of Problem 3}

The $P$ deviation of $10^6$ MC runs is $\le3\times10^{-4}$, and the linearization residual is $0.1293^\circ$ (about $0.2\%$ of the variance); the exact Hermite $s_{\max}=66.84$~m differs from the closed form $66.92$~m by $0.13\%$.

\subsection{Validation of Problem 4}

The 74-scenario accounting audit finds no anomalies; the $0.5$~m margin, zero-displacement tolerances of $0.02\bsl0.05\bsl0.10$~m, replacement with the $10.837$~s robust baseline, and $22.5^\circ$ densification all leave the $\eta/t_{\mathrm{last}}$ conclusions unchanged.

% =====================================================================
% 6. Sensitivity Analysis
% =====================================================================
\section{Sensitivity Analysis}
\label{sec:sensitivity}

\subsection{Parameter Selection and Perturbation Design}

Parameters are selected by relevance to the problem, dimensional scale, model dominance, and interpretability: geometric parameters $H$, $\varepsilon_{\max}$, $R_{\mathrm{veh}}$, and B2 height; accuracy parameters $\sigma_\alpha$, $d$, and $\gamma$; timeliness parameters $v$, $\omega$, and the 30~s limit; disturbance parameters $t_{\mathrm{rec}}$, failed vehicle, new-alert $\theta/\rho$, and safety margin. Continuous parameters use one-factor perturbations, $v\times\omega$ uses a $3\times3$ full factorial, discrete disturbances use the 74-scenario enumeration, and all other parameters stay at nominal values.

\subsection{Main Results}

\begin{table}[htbp]
  \centering
  \caption{Objective time under the $v\times\omega$ full factorial (unit: s)}
  \label{tab:6-1}
  \begin{tabular}{@{}lccc@{}}
    \toprule
    $v\backslash\omega$ & $24^\circ$/s & $30^\circ$/s & $36^\circ$/s \\
    \midrule
    $2.5$~m/s & 11.655 & 10.611 & 9.932 \\
    $3.0$~m/s & 11.499 & 10.355 & 9.676 \\
    $3.5$~m/s & 11.387 & 10.243 & 9.493 \\
    \bottomrule
  \end{tabular}
\end{table}

\noindent\textbf{Problem 1 elasticities.} The closed-form partial derivatives of $r^*$ are $\partial r^*/\partial H=0.1763$~m/m, $\partial r^*/\partial\varepsilon_{\max}=-0.7738$~m/deg, and $\partial r^*/\partial h_{\mathrm{cam}}=-0.1763$~m/m; at the nominal point $(H,\varepsilon_{\max})=(45~\mathrm{m},80^\circ)$, the elasticities are $E_H=1.05$ and $E_{\varepsilon\max}=(-0.7738\times80)/7.5821\approx-8.16$. As $H$ increases from 30 to 60~m, $r^*$ increases linearly from $4.94$ to $10.23$~m; as $\varepsilon_{\max}$ increases from $70^\circ$ to $85^\circ$, $r^*$ decreases nonlinearly from $15.65$ to $3.76$~m. Replacing $Q$ with $Q_2=(1/r)\cos^2(\varepsilon-45^\circ)$ leaves the argmax $(92.2,69.4)$, $r=7.5828$~m, and the optimal set unchanged; on the $\rmin$ circle, $Q_2$ is identically $0.0885$~m$^{-1}$, showing that the criterion choice is robust.

\noindent\textbf{Problem 2 elasticities.} The elasticity of $\omega$ is about $-0.44$ and that of $v$ about $-0.11$, so the former is about four times the latter; a $0.5$~m margin raises the objective from $10.355$ to $10.837$~s ($+4.7\%$) and increases the clearance from $1.3$~mm/$6.7$~cm to more than $0.5$~m. The sensitivity of $G^*$ to $H$ has the same origin as in Problem~1 ($0.1763$~m/m); $\sigma_\alpha$ and DRMS are strictly linear, and $G$ doubles at $\gamma=30^\circ$.

\noindent\textbf{Slope in $d$ for Problem 3.} $d=30\bsl35\bsl40\bsl45\bsl50\bsl60$~m corresponds to $s_{\max}=46.44\bsl56.68\bsl66.92\bsl77.16\bsl87.40\bsl107.88$~m, with a slope of $2.0481$~m/m; confidence levels $0.90\bsl0.95\bsl0.98\bsl0.99$ correspond to $85.20\bsl66.92\bsl51.41\bsl43.07$~m; gimbal-error scalings $0.5\bsl0.75\bsl1.0\bsl1.25\bsl1.5$ correspond to $72.32\bsl70.11\bsl66.92\bsl62.63\bsl57.03$~m. $d$ is the strongest continuous lever, the confidence level acts nonlinearly, and gimbal improvements saturate over long travels.

\subsection{Stability Conclusions}

Within reasonable perturbation ranges, no qualitative conclusion flips: the optimum lies on the boundary, $G^*=\rmin$, localization dominates over long travel, and rescheduling converts structural loss into temporal loss all remain valid. The model is insensitive to numerical/engineering parameters such as the grid step, $0.5$~m margin, zero-displacement tolerance, and $45^\circ$ azimuth sampling; the two-stage response of $t_{\mathrm{last}}$ to $t_{\mathrm{rec}}$ in Problem~4 is shown in Section~4.5.3. Remodeling is needed only under extreme conditions (the $\rmin$ circle is completely occluded, $t_{\mathrm{rec}}$ approaches 30~s with triple disturbances, or errors are strongly correlated).

% =====================================================================
% 7. Strengths and Weaknesses
% =====================================================================
\section{Strengths and Weaknesses}
\label{sec:strengths}

\subsection{Strengths}

\begin{enumerate}[label=(\arabic*)]
\item \textbf{Analytical certificates instead of numerical search.} Problem~1 derives $r^*=7.5821$~m from C2 and proves attainability on the feasible arc; Problem~2(1) proves $G^*=\rmin$ with the QM--AM inequality and C2; Problem~3 solves $s_{\max}$ in closed form from the monotonicity of $\sigma_\psi(s)$ and verifies it by exact quadrature. All three give a complete ``lower bound + attainability/equality'' proof rather than a best grid value.

\item \textbf{Integrated geometry--accuracy--timing--disturbance modeling with cross-validation.} The four problems share the same site geometry, observable region, and observation timing: $\Omega$ is used for station selection, the 2(2) timeline becomes the Problem~4 baseline, and the Problem~4 snapshot reuses the same slew/reset rules. Closed form--grid--sampling, covariance--MC--path simulation, linearization--Hermite--$10^6$ MC, and analytical benchmark--74 scenarios--robust re-solution all have residuals at the discretization/linearization level.

\item \textbf{Robustness and deployability built in.} Problem~2(2) discloses the $1.3$~mm/$6.7$~cm knife-edge clearance and gives a $10.837$~s robust variant; Problem~1 gives a $78^\circ$ fallback design with a $20.5\%$ cost; Problem~3 gives the conservative $66.84$~m value; Problem~4 uses 74 scenarios and $22.5^\circ$ densification for validation. Key numbers are driven by a single JSON data source, and the scripts are reproducible on any machine.
\end{enumerate}

\subsection{Weaknesses}

\begin{enumerate}[label=(\arabic*)]
\item \textbf{The global optimality of Problem~2(2) is not certified.} The assignment structures are exhaustively enumerated at the $10^3$ level, but station optimization relies on coordinate descent and Nelder--Mead; $10.355$~s is the best solution in the sense of ``structural exhaustion + station refinement,'' and the lower bound $3.414$~s is loose because it ignores reuse delays.

\item \textbf{The independent normal error assumption is strong.} Problem~2 assumes independent $\mathcal{N}(0,(1^\circ)^2)$ azimuth errors at the two stations, and Problem~3 assumes three independent zero-mean normal sources; MC validates the nominal parameters but does not cover correlated, nonzero-mean, or heavy-tailed cases.

\item \textbf{The scenario set is a discrete representative set, and the kinematics are simplified.} The 74 scenarios cover 8 azimuths, 3 distances, 5 failure types, and 3 recovery times, but not continuous arrivals, secondary disturbances, or adversarial choices; vehicles travel at constant speed with no acceleration/turning/collision avoidance, the gimbal rotates ideally, and the vehicle--target distance in Problem~3 is fixed at 40~m. These simplifications are consistent with the problem statement but would change the timeline and error propagation of a more realistic system.
\end{enumerate}

\subsection{Model Extensions}

\begin{enumerate}[label=(\arabic*)]
\item \textbf{Multi-station/multi-target.} The two-station covariance $\Sigma=A^{-1}DA^{-T}$ can be extended to multi-station intersection by stacking the rows of $A$; for multiple targets, each target has its own $\Sigma$ and shares the same vehicle and timeline constraints.

\item \textbf{Arbitrary monotone $Q(r)$.} The ``closed-form lower bound + attainability'' certificate of Problem~1 depends only on the monotonicity of $Q$ in $r$; the $Q_2$ comparison in Section~6.2 shows that switching to another monotone criterion in the same family leaves the optimal set unchanged, so the certificate transfers to other monotone quality measures.

\item \textbf{Continuous arrivals/secondary disturbances.} The event-triggered snapshot and re-optimization framework can directly accept continuously arriving new alerts and secondary disturbances; it suffices to refresh $A_p/V_p$ at the trigger time and reuse the same solver.
\end{enumerate}

% =====================================================================
% 8. Conclusion
% =====================================================================
\section{Conclusion}
\label{sec:conclusion}

First, the single-target observable region is determined by the intersection of the vehicle geometry, gimbal pitch, and LOS occlusion constraints; the optimal quality is fixed by the closed-form lower bound $r^*=7.5821$~m ($Q^*=0.1319$~m$^{-1}$), and the optimal set is a continuous arc that must be made unique by tie-breaking.

Second, two-station accuracy is characterized by the geometric dilution of precision $G$, whose lower bound is attained at $d_1=d_2=\rmin$ and $\gamma=90^\circ$; under the 30~s limit, V3+V5 is selected, and the maximum scheduling start is $10.355$~s, $50.6\%$ lower than the nearest-assignment baseline of $20.968$~s; timeliness and accuracy are not interchangeable.

Third, the pointing deviation is composed of localization, bias, and control sources; the maximum travel distance satisfying $P(|\Delta\psi|\le5^\circ)\ge0.95$ is $66.92$~m (with the exact-geometry conservative value $66.84$~m), and localization dominates over long travel.

Fourth, event-triggered rescheduling improves the completion rate in 67 of 74 scenarios, ties in 7, and degrades none; its value lies in converting structural loss into temporal loss.

Overall, geometry determines feasibility, accuracy determines station placement, scheduling determines timeliness, and rescheduling preserves the completion rate.

% =====================================================================
% References
% =====================================================================
\begin{thebibliography}{9}
\bibitem{swmcm} 2026 12th ShuWei Cup National College Mathematical Modeling Challenge (Autumn), Problem B: Rapid UAV Localization via Coordinated Mobile Vehicles. 2026.

\bibitem{preparata} Preparata F P, Shamos M I. Computational Geometry: An Introduction. New York: Springer, 1985.

\bibitem{torrieri} Torrieri D J. Statistical theory of passive location systems. IEEE Transactions on Aerospace and Electronic Systems, 1984, AES-20(2): 183--198.

\bibitem{foy} Foy W H. Position-location solutions by Taylor-series estimation. IEEE Transactions on Aerospace and Electronic Systems, 1976, AES-12(2): 187--194.

\bibitem{dijkstra} Dijkstra E W. A note on two problems in connexion with graphs. Numerische Mathematik, 1959, 1: 269--271.

\bibitem{nelder} Nelder J A, Mead R. A simplex method for function minimization. The Computer Journal, 1965, 7(4): 308--313.

\bibitem{hammersley} Hammersley J M, Handscomb D C. Monte Carlo Methods. London: Methuen, 1964.
\end{thebibliography}


% ---------------- appendix: separate numbering A, B, C ------------
\clearpage
\pagenumbering{Alph}
\rhead{\small Appendix Page \thepage}
% =====================================================================
% Appendix (separate page numbering: A, B, C)
% =====================================================================
\section*{Appendix A: Proofs and Supplementary Derivations}
\addcontentsline{toc}{section}{Appendix A: Proofs and Supplementary Derivations}

\subsection*{A.1 Problem 1: Intersection of the Optimal Arc and the Proof That C3 Adds No Occlusion}
\addcontentsline{toc}{subsection}{A.1 Problem 1: Intersection of the optimal arc and C3}

\noindent\textbf{(a) Derivation of the blocked arc.} Let $T'=(84.6410,70.0000)$ and $\rmin=7.5821$~m. The $1.5$~m expanded rectangle of B2 is $[61.5,80.5]\times[57.5,72.5]$; a point on the circle, $(84.6410+7.5821\cos\varphi,\ 70+7.5821\sin\varphi)$, falls inside this rectangle if and only if
\begin{equation}
  \cos\varphi\le\frac{80.5-84.6410}{7.5821}=-0.54616,
  \qquad
  \sin\varphi\le\frac{72.5-70}{7.5821}=0.32972.
\end{equation}
The first inequality gives $\varphi\in[123.11^\circ,236.89^\circ]$, and the second gives $\varphi\in[0^\circ,19.25^\circ]\cup[160.75^\circ,360^\circ]$ ($x\ge61.5$ and $y\ge57.5$ hold automatically on the circle because $x\ge77.06$ and $y\ge62.42$ there). Their intersection is $\varphi\in[160.75^\circ,236.89^\circ]$, i.e., a blocked arc of $76.15^\circ$. Hence C1 leaves the continuous arc $[236.89^\circ,360^\circ)\cup(0^\circ,160.75^\circ]$ on the circle, of total length $283.85^\circ$.

\noindent\textbf{(b) C3 adds no occlusion on the feasible arc.} Basis: the segment $pT'$ is contained in the disk centered at $T'$ with radius $\rmin$ (convexity of the disk); the distances from the other buildings to $T'$ are $67.09$, $60.64$, $36.00$, and $41.32$~m, all greater than $\rmin+1.5=9.0821$~m (see Section~4.2.3), so only B2 needs to be excluded. For $p$ on the feasible arc, consider two cases.

\noindent (i) $x(p)\ge79$: the $x$-coordinate of the horizontal projection of the LOS lies between $x(p)$ and $84.6410$, hence is always $\ge79$ and outside the east face of B2.

\noindent (ii) $x(p)<79$: by the C1 test in (a), $p$ must satisfy $y(p)>72.5$. Let $p=(x_0,y_0)$, $t=y_0-70\in(2.5,7.5821]$, and $x_0=84.6410-\sqrt{\rmin^2-t^2}$. The ordinate of the intersection of $pT'$ with the line $x=79$ reduces to
\begin{equation}
  y(79)=70+\frac{5.6410\,t}{\sqrt{\rmin^2-t^2}}.
\end{equation}
This expression is strictly increasing on $t\in(2.5,7.5821]$, so $y(79)\ge y(79)|_{t=2.5}=71.97>71$; that is, the LOS passes north of the north face of B2 ($y=71$) and does not enter the B2 footprint. Therefore C3 produces no additional occlusion on the feasible arc.

\subsection*{A.2 Problem 2(1): Algebraic Expansion of $\Sigma=A^{-1}DA^{-T}$}
\addcontentsline{toc}{subsection}{A.2 Problem 2(1): Algebraic expansion of the covariance}

Let the LOS azimuths of stations $p_1$ and $p_2$ be $\alpha_i$, with unit directions $u_i=(\cos\alpha_i,\sin\alpha_i)$ and unit normals $n_i=(-\sin\alpha_i,\cos\alpha_i)$ ($n_i^Tu_i=0$). A first-order azimuth error $\delta\alpha_i$ makes the signed distance from the true target to the perturbed bearing line approximately $d_i\delta\alpha_i$, so
\begin{equation}
  n_i^T\Delta p=d_i\delta\alpha_i,\qquad i=1,2,
\end{equation}
which in matrix form is $A\Delta p=\delta$ with
\begin{equation}
  A=\begin{bmatrix}n_1^T\\ n_2^T\end{bmatrix}
  =\begin{bmatrix}n_{1x}&n_{1y}\\ n_{2x}&n_{2y}\end{bmatrix},
  \qquad
  \delta=(d_1\delta\alpha_1,d_2\delta\alpha_2)^T.
\end{equation}
Since $\delta\alpha_1$ and $\delta\alpha_2$ are independent and identically distributed $\mathcal{N}(0,\sigma^2)$, $\mathbb{E}[\delta\delta^T]=D=\diag(d_1^2\sigma^2,d_2^2\sigma^2)$, and
\begin{equation}
  \Sigma=\mathbb{E}[\Delta p\Delta p^T]
  =A^{-1}\mathbb{E}[\delta\delta^T]A^{-T}=A^{-1}DA^{-T}.
\end{equation}

\noindent\textbf{(a) Inverse and columns.} $\det A=n_{1x}n_{2y}-n_{1y}n_{2x}=n_1\times n_2=\sin\gamma$ ($\gamma$ is the angle between the two LOS; the signed area is $\pm\sin\gamma$, and its absolute value enters the variance). The $2\times2$ adjugate is
\begin{equation}
  \operatorname{adj}A=\begin{bmatrix}n_{2y}&-n_{1y}\\ -n_{2x}&n_{1x}\end{bmatrix},
\end{equation}
so
\begin{equation}
  A^{-1}=\frac{1}{\det A}\begin{bmatrix}n_{2y}&-n_{1y}\\ -n_{2x}&n_{1x}\end{bmatrix}.
\end{equation}
Its two columns are
\begin{equation}
  \mathrm{col}_1(A^{-1})=\frac{1}{\det A}(n_{2y},-n_{2x})^T=\frac{1}{\det A}u_2,
\end{equation}
\begin{equation}
  \mathrm{col}_2(A^{-1})=\frac{1}{\det A}(-n_{1y},n_{1x})^T=-\frac{1}{\det A}u_1,
\end{equation}
using the $90^\circ$ rotation relation between $u_i=(\cos\alpha_i,\sin\alpha_i)$ and $n_i=(-\sin\alpha_i,\cos\alpha_i)$.

\noindent\textbf{(b) Outer-product expansion.}
\begin{equation}
  \Sigma=A^{-1}DA^{-T}
  =\sum_{i=1}^{2}D_{ii}\,\mathrm{col}_i(A^{-1})\,\mathrm{col}_i(A^{-1})^T
  =\frac{\sigma^2}{\det^2A}\left[d_1^2u_2u_2^T+d_2^2u_1u_1^T\right].
\end{equation}
No cross term appears because $D$ is diagonal; if $u_i$ were mistakenly written as the normal $n_i$, the expression would be rotated by $90^\circ$ (trace and determinant unchanged, but the matrix itself different). The main text consistently uses the LOS direction $u_i$.

\noindent\textbf{(c) Trace and determinant.}
\begin{equation}
  \trace\Sigma=\frac{\sigma^2}{\det^2A}\left[d_1^2\|u_2\|^2+d_2^2\|u_1\|^2\right]
  =\frac{\sigma^2(d_1^2+d_2^2)}{\sin^2\gamma},
\end{equation}
\begin{equation}
  \det\Sigma=\frac{\det D}{\det^2A}=\frac{\sigma^4d_1^2d_2^2}{\sin^2\gamma}.
\end{equation}
These correspond respectively to
\begin{equation}
  \mathrm{DRMS}^2=\frac{\trace\Sigma}{2}
  =\frac{\sigma^2(d_1^2+d_2^2)}{2\sin^2\gamma}
  \ \Longrightarrow\
  \mathrm{DRMS}=\frac{\sigma\sqrt{(d_1^2+d_2^2)/2}}{|\sin\gamma|}=\sigma G.
\end{equation}

\noindent\textbf{(d) $1\sigma$ ellipse and $95\%$ radius.} From $\det\Sigma$, the area of the $1\sigma$ error ellipse is
\begin{equation}
  S_{1\sigma}=\pi\sqrt{\det\Sigma}=\frac{\pi\sigma^2d_1d_2}{|\sin\gamma|}.
\end{equation}
When $d_1=d_2=d$, $S_{1\sigma}=\pi\sigma^2dG$, so the two are of the same order; in general, both $G$ and $S_{1\sigma}$ vary monotonically and inversely with $|\sin\gamma|$, so $G$ is a monotone representative of ``area-type'' accuracy indices. The $95\%$ confidence ellipse is
\begin{equation}
  (\Delta p)^T\Sigma^{-1}(\Delta p)\le\chi^2_2(0.95)=5.991,
\end{equation}
with semi-axes $\sqrt{5.991\lambda_i}$, where $\lambda_i$ are the eigenvalues of $\Sigma$. When $\Sigma$ is isotropic ($\lambda_1=\lambda_2=\mathrm{DRMS}^2$), the ellipse degenerates to a circle of radius $2.4477\,\mathrm{DRMS}$; in general, $2.4477\,\mathrm{DRMS}$ is the nominal circle radius, and exact $95\%$ coverage should use the ellipse semi-axes. The main text reports $R_{95}=2.4477\,\mathrm{DRMS}$ and notes ``exact when isotropic.''

\noindent\textbf{(e) Degeneracy and condition number.} As $\gamma\to0^\circ$ or $180^\circ$, $\det A=\pm\sin\gamma\to0$, $A$ becomes singular, and the entries of $A^{-1}$ diverge as $1/|\sin\gamma|$; by (c), $\trace\Sigma$ and $\det\Sigma$ both diverge as $1/\sin^2\gamma$. $\det A=n_1\times n_2=\pm\sin\gamma$ measures the singularity of $A$; using $1/|\det A|=1/|\sin\gamma|$ as the condition-number indicator (same order as the 2-norm condition number $\cot(\gamma/2)$, both diverging as $1/\gamma$ as $\gamma\to0$) gives the quantitative basis for the operating window $[30^\circ,150^\circ]$: at $\gamma=30^\circ$, $G=2G(90^\circ)$, i.e., the accuracy cost doubles.

\subsection*{A.3 Problem 3: Linearization Residual RMS $=\sigma(s)^2/d^2$ and Hermite Quadrature}
\addcontentsline{toc}{subsection}{A.3 Problem 3: Linearization residual and Hermite quadrature}

\noindent\textbf{(a) Exact angle and first-order approximation.} Take the true position as the origin and the true LOS as the $x'$ axis, with the target at $(d,0)$, $d=40$~m. Decompose the localization error into along-range $\delta_\parallel$ and cross-range $\delta_\perp$, isotropic and independent, $\delta_\parallel,\delta_\perp\sim\mathcal{N}(0,\sigma(s)^2)$. The exact azimuth deviation from the estimated position to the target is
\begin{equation}
  \Delta\alpha=\operatorname{atan2}(-\delta_\perp,d-\delta_\parallel).
\end{equation}
Let $u=\delta_\parallel/d$ and $v=\delta_\perp/d$; then $d-\delta_\parallel=d(1-u)$, and
\begin{equation}
  \Delta\alpha=-\arctan\frac{v}{1-u}.
\end{equation}
Expanding $v/(1-u)=v(1+u+u^2+O(u^3))$ for $|u|<1$ and using $\arctan z=z-z^3/3+O(z^5)$,
\begin{equation}
  \Delta\alpha=-\left[v+uv+u^2v-\frac{v^3}{3}+O(\|(u,v)\|^4)\right].
\end{equation}
The first-order (linearized) model takes $\Delta\alpha\approx-v=-\delta_\perp/d$.

\noindent\textbf{(b) Residual and RMS.} The residual is
\begin{equation}
  R=\Delta\alpha-(-v)=-uv-u^2v+\frac{v^3}{3}+O(\|(u,v)\|^4).
\end{equation}
The leading term is $-uv$. Since $\delta_\parallel$ and $\delta_\perp$ are independent, zero-mean, and have variance $\sigma^2$,
\begin{equation}
  \mathbb{E}[R^2]=\mathbb{E}[u^2v^2]+O(\sigma^6/d^6)
  =\frac{\mathbb{E}[\delta_\parallel^2]\,\mathbb{E}[\delta_\perp^2]}{d^4}+O(\sigma^6/d^6)
  =\frac{\sigma^4}{d^4}+O(\sigma^6/d^6).
\end{equation}
Hence the residual RMS (in radians) is
\begin{equation}
  \mathrm{RMS}_R=\sqrt{\mathbb{E}[R^2]}=\frac{\sigma(s)^2}{d^2}+O(\sigma(s)^3/d^3).
\end{equation}
Note that $\sigma^2/d^2$ is dimensionless; under the small-angle approximation its numerical value is the RMS in radians, and conversion to degrees multiplies by $180/\pi$. At $s=80$~m, $\sigma=1.9$~m, and $d=40$~m:
\begin{equation}
  \mathrm{RMS}_R=\left(\frac{1.9}{40}\right)^2=2.256\times10^{-3}~\mathrm{rad}=0.1293^\circ.
\end{equation}
Relative to $\sigma_\psi(80)=2.8995^\circ$ at the same point, the standard-deviation share is $0.1293/2.8995=4.46\%$, and the variance share is $(0.1293/2.8995)^2=0.199\%\approx0.2\%$. Higher-order terms are $O(\sigma^3/d^3)$, and at $\sigma/d\le1.9/40=0.0475$ they are at most about $5\%$ of the leading term, so the first-order residual estimate is consistent within the parameter range of this paper.

\noindent\textbf{(c) Normal composition and the $P$ criterion.} The gimbal bias $\varepsilon_0\sim\mathcal{N}(0,0.8^2)$, the control error $\varepsilon_c\sim\mathcal{N}(0,0.6^2)$, and the localization error are independent, so $S=\varepsilon_0+\varepsilon_c\sim\mathcal{N}(0,1.0^2)$ (in degrees). Given $(\delta_\parallel,\delta_\perp)$, $\Delta\alpha$ is deterministic, and the conditional probability is
\begin{equation}
  P(|\Delta\alpha+S|\le5^\circ\mid\delta_\parallel,\delta_\perp)
  =\Phi\!\left(\frac{5^\circ-\Delta\alpha}{1^\circ}\right)
  -\Phi\!\left(\frac{-5^\circ-\Delta\alpha}{1^\circ}\right),
\end{equation}
where $\Phi$ is the standard normal CDF. Therefore
\begin{equation}
  P(s)=\mathbb{E}_{\delta\parallel,\delta\perp}\!\left[
  \Phi\!\left(\frac{5^\circ-\Delta\alpha(\delta_\parallel,\delta_\perp)}{1^\circ}\right)
  -\Phi\!\left(\frac{-5^\circ-\Delta\alpha(\delta_\parallel,\delta_\perp)}{1^\circ}\right)\right].
\end{equation}
Since $\delta_\parallel$ and $\delta_\perp$ are independent and have equal variance, the two-dimensional expectation can be evaluated by tensor-product Gauss--Hermite quadrature:
\begin{equation}
  P(s)\approx\frac{1}{\pi}\sum_{i=1}^{256}\sum_{j=1}^{256}w_iw_j
  \left[\Phi\!\left(\frac{5^\circ-\Delta\alpha(\sqrt2\sigma(s)x_i,\sqrt2\sigma(s)x_j)}{1^\circ}\right)
  -\Phi\!\left(\frac{-5^\circ-\Delta\alpha(\cdot)}{1^\circ}\right)\right],
\end{equation}
where $x_i$ and $w_i$ are the 256-point probabilistic Hermite nodes and weights ($\int e^{-x^2}f(x)\,dx\approx\sum w_if(x_i)$, $\sum w_i=\sqrt\pi$). $P(s)$ decreases monotonically with $s$ ($\sigma(s)=a+bs$ is increasing, and the distribution of $\Delta\alpha$ widens with $\sigma$); bisection/root finding on $P(s)=0.95$ gives $s_{\max}=66.84$~m. The closed-form linearization gives $66.92$~m, a $0.13\%$ difference consistent with the Monte Carlo results in Table~\ref{tab:4-5}.

% ---------------------------------------------------------------------
\section*{Appendix B: Algorithm Routes and Main Program List}
\addcontentsline{toc}{section}{Appendix B: Algorithm Routes and Main Program List}

\subsection*{B.1 Algorithm routes (full table)}
\renewcommand{\thetable}{B-\arabic{table}}
\setcounter{table}{0}
\begin{table}[htbp]
\centering
\caption{Algorithm routes and selection rationale for the four problems (full)}
\label{tab:B-1}
{\footnotesize
\renewcommand{\arraystretch}{1.08}
\begin{tabularnb}{@{}p{10mm}p{43mm}p{37mm}p{26mm}p{32mm}@{}}
  \toprule
  Problem & Algorithm (key steps) & Rationale & Scale / complexity & Comparison with alternatives \\
  \midrule
  Q1 & Closed-form lower bound (C2 $\Rightarrow r\ge\rmin$); 200,000-point sampling on the $\rmin$ circle; 3-D slab intersection; visibility graph + Dijkstra; three-level lexicographic rule &
  $Q=1/r$ is strictly monotone, so the optimal value is fixed in closed form by the active constraint; the optimal set is a one-dimensional arc, and sampling only describes the set, not a 2-D search &
  Closed form $O(1)$; arc sampling $O(2\times10^5)$; one grid verification $O(10^6)$ &
  A 0.1~m grid only gives the best grid value 7.5828~m; the closed form also gives the global certificate 7.5821~m \\
  \midrule
  Q2(1) & Covariance propagation $\Sigma=A^{-1}DA^{-T}$; QM--AM lower-bound certificate; enumeration of $C(5,2)=10$ pairs; polar candidate stations + Nelder--Mead &
  The equality condition reduces optimality verification to a one-dimensional $90^\circ$ pairing test on the $\rmin$ circle; with only five vehicles, pairs can be exhaustively enumerated &
  10 pairs; 4-D NM $\le300$ iterations per pair &
  Direct 4-D search has no optimality proof and easily falls into local optima; the certificate gives $G^*=5.8188$~m \\
  \midrule
  Q2(2) & 1,000 assignment structures expanded through reuse-chain orderings into 5,880 combinations; fixed-point timeline simulation; greedy initialization $\to$ Gauss--Seidel coordinate descent $\to$ Nelder--Mead; top-12 re-refinement; deadlock detection &
  With only 5 vehicles and 6 slots, complete enumeration exhausts the structures and is deterministic and reproducible; continuous station positions are handled by local refinement &
  5,880 simulations; at most 2 tasks per vehicle; $\le160$ NM iterations per call &
  10.355~s vs.\ 20.968~s for nearest assignment ($-50.6\%$); certified LB 7.391~s (gap 28.6\%); randomized searches such as GA cannot guarantee the same solution \\
  \midrule
  Q3 & Closed-form first-order propagation; two-sided quantile criterion; bisection; $256\times256$ Gauss--Hermite quadrature; $10^6$ MC verification &
  $\sigma_\psi(s)$ is monotone, so the $P$ criterion can be inverted in closed form; Hermite quadrature has no sampling error and uses $1/15$ of the MC nodes &
  Closed form $O(1)$; 65,536 quadrature nodes; $10^6$ MC for verification only &
  Closed form 66.92~m vs.\ exact 66.84~m (0.13\%); direct MC root search requires multiple rounds of $10^6$ samples \\
  \midrule
  Q4 & Event-triggered snapshot; reuse of the Q2(2) solver; deterministic enumeration of 74 scenarios; keep/simple-repair/re-optimization comparison &
  After a disturbance, available resources and pending alerts are of the same order as Q2(2); reusing the solver keeps the conventions consistent; the scenario set is small, so no online stochastic optimization is needed; simple repair separates reaction from re-optimization &
  74 scenarios $\times$ one re-optimization each; timeline audit finds 0 negative durations/overlaps &
  Simple repair restores $\eta$ in 67/74; re-optimization lowers $t_{\mathrm{last}}$ in 61/74 (mean 4.35~s) \\
  \bottomrule
\end{tabularnb}
}
\end{table}

\subsection*{B.2 Main program list}
\begin{itemize}
\item \textbf{Problem 1} (\texttt{Q1/}): \texttt{q1\_solve.py} (domain gridding, arc sampling, tie-breaking, sensitivity); \texttt{area\_estimators.py}, \texttt{q1\_metric\_compare.py}, \texttt{export\_params.py} (three-convention area recomputation, $Q/Q_2$ comparison, export of the B1--B5/V1--V5/T1--T3 parameters).
\item \textbf{Problem 2} (\texttt{Q2/}): \texttt{q2a\_core.py}, \texttt{q2b\_core.py}, \texttt{q2b\_gamma.py}, \texttt{q2b\_bound.py} (two-station accuracy, vehicle and station selection, six-slot five-vehicle scheduling and baseline, 2(2) intersection-geometry check, reuse-aware certified lower bound).
\item \textbf{Problem 3} (\texttt{Q3/}): \texttt{q3\_solve.py} (error propagation, $s_{\max}$, the $(a,b)$ region, Monte Carlo, quadrature).
\item \textbf{Problem 4} (\texttt{Q4/}): \texttt{q4\_core.py}, \texttt{q4\_repair.py}, \texttt{q4\_robust.py}, \texttt{q4\_tlast\_diff.py}, \texttt{q4\_scenario\_list.py} (state snapshot, rescheduling, simple-repair baseline, 74 scenarios, robustness checks, timeline comparison, scenario details).
\item \textbf{Build scripts}: \texttt{build\_numbers.py} $\to$ \texttt{numbers.md} (single data-source mapping from body numbers to JSON fields and scripts); \texttt{params.csv} / \texttt{params\_table.md} (parameter specification exported by \texttt{export\_params.py} and cross-checked with \texttt{q1\_solve.py} and \texttt{q2\_lib.py}).
\end{itemize}

\noindent The directory names \texttt{Q1/}--\texttt{Q4/} stand for the Chinese original directories in the submitted supporting materials. The correspondence between body numbers and data files is as follows:
\begin{itemize}
\item Section~4.2.4, three area conventions $\to$ \texttt{Q1/results/q1\_area\_estimators.json} (\texttt{area\_estimators.py}).
\item Table~\ref{tab:4-4}, the clearance-robust variant geometry of Section~4.3.5, and the $95\%$ ellipse semi-axes $\to$ \texttt{Q2/results/q2b\_precision.json} (\texttt{q2b\_gamma.py}).
\item Section~4.3.7, the certified lower bound and gap $\to$ \texttt{Q2/results/q2b\_lower\_bound.json} (\texttt{q2b\_bound.py}).
\item Section~4.4.3 and Table~\ref{tab:4-5} $\to$ \texttt{Q3/results/q3\_summary.txt} (\texttt{q3\_solve.py}).
\item Section~4.5.3 and Table~\ref{tab:4-6}, scenario details $\to$ \texttt{Q4/results/q4\_scenario\_list.md} (\texttt{q4\_scenario\_list.py}).
\item Section~4.5.3, simple-repair baseline and the three-way comparison $\to$ \texttt{Q4/results/q4\_repair.json} (\texttt{q4\_repair.py}).
\item Section~6.2, $Q/Q_2$ comparison $\to$ \texttt{Q1/results/q1\_metric\_compare.json} (\texttt{q1\_metric\_compare.py}).
\item Tables D-1 and D-2, all ten pairs and the 2(2) timeline $\to$ \texttt{Q2/results/q2a\_results.json}, \texttt{Q2/results/q2b\_results.json}.
\end{itemize}

\noindent The complete source programs, run scripts, and result JSON files are provided in the supporting materials.

\subsection*{B.3 Full symbol table}
\begin{table}[htbp]
\centering
\caption{Full symbol list}
\label{tab:B-2}
{\small
\renewcommand{\arraystretch}{1.08}
\begin{tabularnb}{@{}p{38mm}p{92mm}p{30mm}@{}}
  \toprule
  Symbol & Meaning & Unit / value \\
  \midrule
  $O$ & Dispatch center & $(50,50)$~m \\
  $B_i$ & Building $i$ ($i=1,\dots,5$), planar range $[x_i^0,x_i^1]\times[y_i^0,y_i^1]$, height $h_i$ & m \\
  $T_a$; $\theta_a,R_a,H_a$; $T'_a$ & Alert $a$ and its target; azimuth (counterclockwise from east), horizontal distance, height; horizontal projection & $^\circ$, m, m; m \\
  $h_{\mathrm{cam}}$; $\Delta z_a$ & Camera height above ground; camera--target height difference $\Delta z_a=H_a-h_{\mathrm{cam}}$ & 2~m; m \\
  $p=(x,y)$; $r(p)$; $\varepsilon(p)$ & Center of the observation station; horizontal distance from camera to target projection; LOS elevation $\arctan(\Delta z_a/r)$ & m; m; $^\circ$ \\
  $[\varepsilon_{\min},\varepsilon_{\max}]$ & Feasible gimbal pitch range & $[-10^\circ,80^\circ]$ \\
  C1, C2, C3; $\Omega_a$ & Vehicle geometry, gimbal pitch, and LOS occlusion constraints; observable region $\Omega_a=\mathrm{C1}\cap\mathrm{C2}\cap\mathrm{C3}$ & ---; m$^2$ \\
  $Q(p)$; $p^*,r^*,\varepsilon^*,Q^*$ & Observation quality $1/r$; optimal station, horizontal distance, elevation, and quality in Problem~1 & m$^{-1}$; m, m, $^\circ$, m$^{-1}$ \\
  $m(p)$; $\varphi$ & Geometric safety margin; polar angle parameter of the optimal arc & m; $^\circ$ \\
  $V_j$; $V_j^0$ & Observation vehicle $j$ and its initial position ($j=1,\dots,5$) & m \\
  $R_{\mathrm{veh}}$; $v$; $\omega$ & Equivalent vehicle disk radius; travel speed; maximum gimbal single-axis angular rate & 1.5~m; 3~m/s; $30^\circ$/s \\
  $\alpha$ & Azimuth from station to target (pan command angle) & $^\circ$ \\
  $t_{\mathrm{tr}},t_{\mathrm{sl}},t_{\mathrm{obs}},t_{\mathrm{rs}}$ & Travel, slew, synchronized observation, and reset times & s ($t_{\mathrm{obs}}=5$) \\
  $d_i$; $d$ & Horizontal distance from station $i$ to target; single-station shorthand in Problem~3 (nominal 40~m) & m \\
  $\sigma_\alpha$; $\gamma$; $\beta$ & Single-station azimuth standard deviation; intersection angle of the two LOS; two-station baseline length & $1^\circ$ (rad in computation); $^\circ$; m \\
  $\Sigma$; $G$ & Covariance of the intersection-point position error; geometric dilution of precision $\sqrt{(d_1^2+d_2^2)/2}/|\sin\gamma|$ & m$^2$; m \\
  DRMS; $R_{95}$ & Root-mean-square position error $\sigma_\alpha G$; 95\% error-circle radius $2.4477\,\mathrm{DRMS}$ & m; m \\
  $\pi_j$; $p_{j,a}$; $M_a$ & Task order of vehicle $j$; station of vehicle $j$ for alert $a$; vehicle pair observing alert $a$ & ---; m; --- \\
  $t^{\mathrm{rd}}_{v,k}$; $t_a^{\mathrm{start}}$ & Ready time of the $k$-th task of vehicle $v$; observation start time of alert $a$ & s \\
  $P,\hat P,\delta$; $s$; $\sigma(s)=a+bs$ & True position, estimated position, and localization error; cumulative travel distance; single-axis localization standard deviation ($a,b\ge0$) & m; m; m, m/m \\
  $\varepsilon_0,\varepsilon_c$; $\Delta\psi$; $\sigma_\psi$; $s_{\max}$ & Gimbal bias/control error; difference between commanded and true azimuth ($\wrap$); combined standard deviation of $\Delta\psi$; maximum travel distance satisfying $P(|\Delta\psi|\le5^\circ)\ge0.95$ & $0.8^\circ/0.6^\circ$; $^\circ$; $^\circ$; m \\
  $t_p,t_{\mathrm{rec}}$; $N,k,\eta$; $t_{\mathrm{last}}$ & Rescheduling trigger/communication recovery time; total alerts, successful starts, $\eta=k/N$; latest observation start time (60~s if not started) & s; ---; s \\
  \bottomrule
\end{tabularnb}
}
\end{table}

% ---------------------------------------------------------------------
\section*{Appendix C: Figure and Table Inventory}
\addcontentsline{toc}{section}{Appendix C: Figure and Table Inventory}

The four problem figure sets contain 11 multi-panel combined figures (3 for Q1, 4 for Q2, 2 for Q3, and 2 for Q4), plus one framework figure in Chapter~2 (Figure~\ref{fig:framework}); each is delivered as an SVG vector with a PNG fallback.

\renewcommand{\thetable}{C-\arabic{table}}
\setcounter{table}{0}

\begin{table}[htbp]
\centering
\caption{Figure inventory}
\label{tab:C-1}
{\small
\renewcommand{\arraystretch}{1.08}
\begin{tabularnb}{@{}p{22mm}p{110mm}p{20mm}@{}}
  \toprule
  Figure & Content & Section \\
  \midrule
  Figure 2-1 & Framework of the four problems and the modeling chain & 2 \\
  Figure 4-1 & Observable region of T1 and building occlusion shadows & 4.2.4 \\
  Figure 4-2 & Optimal arc and margin--azimuth curve & 4.2.3 \\
  Figure 4-3 & Sensitivity of $r^*$ and the observable area & 4.2.4 \\
  Figure 4-4 & Observable region of T2 and 30~s reachable domains & 4.3.2 \\
  Figure 4-5 & Two-station accuracy versus $\gamma$, $d$, and baseline & 4.3.2 \\
  Figure 4-6 & Intersection error ellipses and the $95\%$ circle & 4.3.2 \\
  Figure 4-7 & Gantt chart and deployment routes of the time-optimal plan & 4.3.5 \\
  Figure 4-8 & Variance components and shares of $\sigma_\psi$ & 4.4.3 \\
  Figure 4-9 & Probability curve $P(|\Delta\psi|\le5^\circ)$ and the $(a,b)$ region & 4.4.3 \\
  Figure 4-10 & Two-stage response of $t_{\mathrm{last}}$ to $t_{\mathrm{rec}}$ & 4.5.3 \\
  Figure 4-11 & Rescheduling Gantt chart and the 74-scenario comparison & 4.5.3 \\
  \bottomrule
\end{tabularnb}
}
\end{table}

\begin{table}[htbp]
\centering
\caption{Table inventory}
\label{tab:C-2}
{\small
\renewcommand{\arraystretch}{1.08}
\begin{tabularnb}{@{}p{22mm}p{110mm}p{20mm}@{}}
  \toprule
  Table & Content & Section \\
  \midrule
  Table 2-1 & Building geometry and initial vehicle positions & 2.1 \\
  Table 2-2 & Compact algorithm routes (full table in B-1) & 2.6 \\
  Table 3-1 & Core symbols (full table in B-2) & 3.1 \\
  Table 4-1 & Main results of Problem~1 & 4.2.4 \\
  Table 4-2 & Two-station accuracy and station selection & 4.3.2 \\
  Table 4-3 & Scheduling plan versus the nearest-assignment baseline & 4.3.4 \\
  Table 4-4 & Intersection geometry of the timeliness plan & 4.3.6 \\
  Table 4-5 & Monte Carlo verification for Problem~3 & 4.4.3 \\
  Table 4-6 & Three-way comparison by scenario class for Problem~4 & 4.5.3 \\
  Table 4-7 & Comparison of representative scenarios & 4.5.3 \\
  Table 5-1 & Summary of validation for the four problems & 5 \\
  Table 6-1 & Objective time under the $v\times\omega$ full factorial & 6.2 \\
  Table B-1 & Full algorithm routes and selection rationale & Appendix B.1 \\
  Table B-2 & Full symbol list & Appendix B.3 \\
  Table D-1 & All ten vehicle pairs for T2 & Appendix D \\
  Table D-2 & Stations and per-task timeline of the 10.355~s plan & Appendix D \\
  \bottomrule
\end{tabularnb}
}
\end{table}

% ---------------------------------------------------------------------
\section*{Appendix D: Supplementary Tables}
\addcontentsline{toc}{section}{Appendix D: Supplementary Tables}

\renewcommand{\thetable}{D-\arabic{table}}
\setcounter{table}{0}

\begin{table}[htbp]
\centering
\caption{All ten vehicle pairs for T2 in 2(1)}
\label{tab:D-1}
{\footnotesize
\renewcommand{\arraystretch}{1.06}
\begin{tabularnb}{@{}p{14mm}p{26mm}p{14mm}p{16mm}p{16mm}p{24mm}p{26mm}@{}}
  \toprule
  Pair & $d_1/d_2$ (m) & $\gamma$ ($^\circ$) & $G$ (m) & DRMS (m) & $t_1/t_2$ (s) & Feasible within 30~s \\
  \midrule
  V1+V2 & 5.8188 / 11.2688 & 90.000 & 8.9678 & 0.1565 & 19.19 / 30.29 & No \\
  V1+V3 & 5.8188 / 5.8188 & 90.000 & 5.8188 & 0.1016 & 18.71 / 3.49 & Yes \\
  V1+V4 & 5.8188 / 5.8188 & 90.000 & 5.8188 & 0.1016 & 18.29 / 27.64 & Yes \\
  V1+V5 & 5.8188 / 5.8188 & 90.000 & 5.8188 & 0.1016 & 19.47 / 14.17 & Yes \\
  V2+V3 & 5.8188 / 11.2688 & 90.000 & 8.9678 & 0.1565 & 34.75 / 10.08 & No \\
  V2+V4 & 5.8188 / 11.2688 & 90.000 & 8.9678 & 0.1565 & 34.75 / 25.37 & No \\
  V2+V5 & 5.8188 / 11.2688 & 90.000 & 8.9678 & 0.1565 & 34.75 / 16.54 & No \\
  V3+V4 & 5.8188 / 5.8188 & 90.000 & 5.8188 & 0.1016 & 6.49 / 27.64 & Yes \\
  \textbf{V3+V5} & \textbf{5.8188 / 5.8188} & \textbf{90.000} & \textbf{5.8188} & \textbf{0.1016} & \textbf{5.11 / 13.61} & \textbf{Yes (selected)} \\
  V4+V5 & 5.8188 / 5.8188 & 90.000 & 5.8188 & 0.1016 & 28.95 / 15.42 & Yes \\
  \bottomrule
\end{tabularnb}
}
\par\smallskip
\begin{minipage}{\linewidth}
\footnotesize Note: $d_1,d_2$ are the station--target distances, $G$ and DRMS follow Section~4.3.1 with $\sigma=1^\circ=0.017453$~rad, and $t_1,t_2$ are the obstacle-avoiding travel times. The four V2 pairs share the same candidate station, hence their identical $G$. ``Feasible'' means both ready times are within the 30~s limit; among the six feasible pairs, V3+V5 has the smallest total travel time ($18.72$~s).
\end{minipage}
\end{table}

\begin{table}[htbp]
\centering
\caption{Stations and per-task timeline of the 10.355~s plan in 2(2)}
\label{tab:D-2}
{\footnotesize
\renewcommand{\arraystretch}{1.06}
\begin{tabularnb}{@{}p{12mm}p{10mm}p{28mm}p{13mm}p{13mm}p{12mm}p{16mm}p{14mm}p{14mm}@{}}
  \toprule
  Vehicle & Task & Station $(x,y)$ (m) & $\tau_{\mathrm{tr}}$ (s) & $\tau_{\mathrm{sl}}$ (s) & Start (s) & Observation end (s) & Reset end (s) & Clearance (m) \\
  \midrule
  V3 & T1 & (15.2004, 87.2331) & 0.57 & 1.03 & 1.607 & 6.607 & 7.641 & 5.581 \\
  V5 & T1 & (49.9000, 92.8000) & 0.07 & 1.53 & 1.607 & 6.607 & 8.140 & 7.080 \\
  V1 & T2 & (14.8000, 14.3000) & 0.12 & 3.12 & 6.994 & 11.994 & 15.117 & 7.295 \\
  V4 & T2 & (86.9000, 70.9000) & 1.03 & 5.96 & 6.994 & 11.994 & 17.953 & 0.067 \\
  V2 & T3 & (88.1000, 9.5000) & 0.17 & 5.75 & 10.355 & 15.355 & 21.100 & 10.897 \\
  V3 & T3 & (16.7656, 88.6590) & 0.71 & 2.01 & 10.355 & 15.355 & 17.363 & 0.001 \\
  \bottomrule
\end{tabularnb}
}
\par\smallskip
\begin{minipage}{\linewidth}
\footnotesize Note: the observation lasts 5~s and the reset time is the slew back to the east-facing attitude. The start time of an alert is the later of its two vehicles' ready times, so V1 waits from $3.24$~s to $6.994$~s for the longer slew of V4. V3 serves T1 and then T3 (the only reuse chain); T1 is already under way at $t_p=5$~s and is therefore locked in Problem~4. Clearance is the minimum distance from the station to the $1.5$~m-expanded building footprints; the two knife-edge values ($0.001$~m and $0.067$~m) are those discussed in Section~4.3.5.
\end{minipage}
\end{table}


\end{document}
